\RequirePackage{fix-cm}
\documentclass[12pt,a4paper]{article}


\usepackage{pdfpages}
\usepackage{graphicx}
\usepackage{tikz} 


\usepackage{fontspec}
\usepackage{polyglossia}

\setmainlanguage[numerals=arabic]{armenian}


\setmainfont{GHEA Grapalat}
\newfontfamily\armenianfont{GHEA Grapalat}
\newfontfamily\armenianfonttt{GHEA Grapalat}


\usepackage[a4paper,
            top=2.5cm,
            bottom=2.5cm,
            left=2.5cm,
            right=2.5cm]{geometry}

\usepackage{setspace}
\usepackage{booktabs}
\usepackage{microtype}
\usepackage{titlesec}
\usepackage{amsmath}
\usepackage{amsfonts}
\usepackage{amssymb}

\numberwithin{equation}{section}
\numberwithin{figure}{section}


\usepackage[hidelinks]{hyperref}
\usepackage[nottoc]{tocbibind}


\titleformat{\section}
{\centering\bfseries\fontsize{14}{16}\selectfont}
{\thesection}{1em}{}

\titleformat{\subsection}
{\bfseries\fontsize{13}{15}\selectfont}
{\thesubsection}{1em}{}


\setlength{\emergencystretch}{3em}
\setlength{\parindent}{0.75cm}
\setlength{\parskip}{0pt}
\setstretch{1.5}
\sloppy


\usepackage{listings}
\usepackage{xcolor}   

\lstdefinelanguage{Julia}
  {morekeywords={struct,end,function,if,else,for,while,using,module,const,let},
   sensitive=true,
   morecomment=[l]{\#},
   morestring=[b]",
  }[keywords,comments,strings]

\lstset{
    backgroundcolor=\color{gray!5},   
    basicstyle=\ttfamily\small,       
    keywordstyle=\color{blue},        
    commentstyle=\color{gray},        
    stringstyle=\color{purple},       
    showstringspaces=false,
    breaklines=true,                  
}


\usepackage{enumitem}
\setlist{nosep, leftmargin=1.25cm}

\begin{document}


\renewcommand{\contentsname}{ԲՈՎԱՆԴԱԿՈՒԹՅՈՒՆ}
\tableofcontents
\newpage

\section*{ՆԵՐԱԾՈՒԹՅՈՒՆ}
\addcontentsline{toc}{section}{ՆԵՐԱԾՈՒԹՅՈՒՆ}

Ներկայացված աշխատանքում ուսումնասիրվում է Ֆերմա--Տորիչելիի խնդիրը: Դրա դասական դրվածքը ձևակերպվել է Ֆերմայի կողմից դեռևս 17-րդ դարում՝ գտնել էվկլիդյան հարթության վրա տրված երեք  կետերից ունեցած հեռավորությունների գումարը մինիմիզացնող կետը: Նույն ժամանակ էլ այդ խնդիրը լուծվել է Տորիչելիի կողմից ([15]): Ավելի ուշ՝ 20-րդ դարում, դիտարկվեց տարբեր մետրիկական տարածություններում դրա բնական ընդհանրացումը  կետերի ավելի մեծ քանակի համար: Ամենամեծ առաջընթացը խնդիրը ստացավ նորմավորված տարածություններում, որոնցում արտածվեցին դրա լուծման արդյունավետ մեթոդներ: Պարզվեց, որ լուծումների բազմությունն իրենից ներկայացնում է տրված կետերից ելնող կոնների հատում և, հետևաբար, կարող է լինել ոչ միակը:

\begin{figure}[htbp]
\centering
\includegraphics[width=0.6\textwidth]{images/1.png}
\caption{Տորիչելիի կետի կառուցումը էվկլիդյան հարթության վրա}
\label{fig:torricelli}
\end{figure}


Ֆերմա--Տորիչելիի պրոբլեմի ուսումնասիրությունն սկսվել է [13] աշխատանքի հետ ծանոթությունից, որում նկարագրված է խնդրի նկատմամբ երկրաչափական մոտեցման կիրառումը և շարադրված են որոշ նոր արդյունքներ, որոնք ստացվել են իրական վերջավորաչափ նորմավորված տարածությունների շրջանակներում: Տվյալ աշխատանքն սկսվում է այդ նյութի համառոտ շարադրանքով: Հիմնական բովանդակությունը բաղկացած է երկու մասից: Այս բնագավառում առաջին խնդիրը դարձավ նորմի վրա այնպիսի պայմանների որոնումը, որոնց դեպքում հանդիպում են ոչ միակ լուծումով կետերի հավաքածուներ: Հիմնական արդյունքն են հանդիսանում նորմավորված տարածությունում ոչ միակ լուծումների գոյության պայմանները: Ապացուցված է չափանիշ՝ տրված կետերի ֆիքսված քանակով խնդրի և կամայական չափայնության տարածության համար: Նաև ճշգրտված են պայմանները հարթությունների և եռաչափ տարածությունների համար, ինչին կցված է չափանիշի կիրառումը որոշ կոնկրետ օրինակների վրա: Այնուհետև  առավել մանրամասն հետազոտվել է ոչ միակ լուծումների տեսքը նորմավորված հարթություններում: Դրվել է այդպիսի լուծումների կայունության հարցը՝ ոչ միակության պահպանումը հավաքածուի կետերի փոքր փոփոխությունների դեպքում: Տվյալ հարցի լուծման համար  սահմանվել է բիֆուրկացիոն դիագրամներ՝ նորմավորված հարթության վրա երեք կետերից բաղկացած հավաքածուների լուծումների համար: Դրանք թույլ են տալիս ակնառու կերպով ցուցադրել ոչ միակ լուծումների կայունությունը համապատասխան նորմավորված հարթություններում: Կայունության հետազոտության արդյունքը դարձավ հարթության վրա նորմի համար չափանիշը: Ստացվել է նաև դասակարգում վեցանկյուն հարթությունների համար, որը լիովին նկարագրում է երեք կետերից բաղկացած հավաքածուների լուծումների տեսքը:
\section{Ֆերմա--Տորիչելիի պրոբլեմը և դրա լուծման մեթոդները}

Այս բաժնի բոլոր պնդումները կձևակերպվեն վերջավորաչափանի նորմավորված տարածությունների համար:

\textbf{Սահմանում 2.1.} $x_0$ կետը կոչվում է \textit{Ֆերմա--Տորիչելիի կետ} $A = \{x_1, \dots, x_n\}$ կետերի համար, եթե $x = x_0$-ն մինիմիզացնում է $\sum_{i=1}^n |x x_i|$-ը: Այդպիսի բոլոր կետերի բազմությունը նշանակենք $\operatorname{ft}(A)$:

$\sum_{i=1}^n |x x_i|$ ֆունկցիայի հատկություններից կարելի է ստանալ հետևյալ պնդումը Ֆերմա--Տորիչելիի խնդրի լուծումների բազմության մասին, տե՛ս, օրինակ, [5]:

\textbf{Պնդում 2.2.} \textit{Դիցուք $A = \{x_1, \dots, x_n\}$-ը կետեր են տարածության մեջ : Այդ դեպքում $\operatorname{ft}(A)$-ն ոչ դատարկ, կոմպակտ և ուռուցիկ բազմություն է:}

Վերջավոր չափանի նորմավորված տարածությունում բազմությունը կոմպակտ է, եթե այն փակ է և սահմանափակ:Ֆերմա-Տորիչելիի խնդրի համատեքստում սա նշանակում է, որ հեռավորությունների գումարը նվազագույնի հասցնող կետերի բազմությունը տարածության մեջ «չի ցրվում» անվերջության և ունի հստակ, վերջավոր չափեր ու սահմաններ:  Բազմությունը ուռուցիկ է, եթե դրա ցանկացած երկու կետերն միացնող հատվածն ամբողջությամբ պատկանում է այդ բազմությանը:  Ֆերմա - Տորիչելիի խնդրի դեպքում սա նշանակում է, որ եթե գոյություն ունեն երկու կամ ավելի օպտիմալ լուծումներ (կետեր), որոնք նվազագույնի են հասցնում հեռավորությունների գումարը, ապա այդ կետերը միացնող ուղղի ցանկացած միջանկյալ կետ նույնպես հանդիսանում է խնդրի վավեր լուծում։

Այժմ կշարադրվի Ֆերմա--Տորիչելիի խնդրի լուծման կառուցման երկրաչափական մեթոդը: Այն օգտագործելու համար, ելակետային $X$ տարածության հետ մեկտեղ, անհրաժեշտ է դիտարկել դրան երկակի $X^*$ տարածությունը, որը բաղկացած է գծային ֆունկցիոնալներից:

\textbf{Սահմանում 2.3.} $\varphi \in X^*$ ֆունկցիոնալը կոչվում է \textit{նորմավորող} $x \in X$ վեկտորի համար, եթե $\|\varphi\| = 1$ և $\varphi(x) = \|x\|$:

$\varphi: X \to \mathbb{R}$ ֆունկցիոնալի նորմը սահմանվում է՝

$$\|\varphi\| = \sup_{\substack{x \in X \\ x \neq 0}} \frac{|\varphi(x)|}{\|x\|} = \sup_{\|x\| \leq 1} |\varphi(x)|$$

Հան-Բանախի թեորեմի հետևանքի համաձայն, ցանկացած ոչ զրոյական $x \in X$ վեկտորի համար միշտ գոյություն ունի առնվազն մեկ այդպիսի $\varphi \in X^*$ նորմավորող ֆունկցիոնալ։

Նշենք նաև, որ ֆունկցիոնալի նորմի համար ճիշտ են նաև հետևյալ առնչությունները՝

$$\|\varphi\| = \sup_{\|x\|=1} |\varphi(x)| = \sup_{\|x\| \leq 1} \varphi(x) = \sup_{\|x\|=1} \varphi(x)$$
Հեշտ է նկատել, որ $X$ տարածության տարրերը, որոնք ստացվել են $x$ վեկտորը դրական $k$ թվով բազմապատկելիս, ունեն նորմավորող ֆունկցիոնալների նույն հավաքածուն: Այսինքն՝ այդպիսի ֆունկցիոնալների բազմությունը կարելի է նկարագրել միավոր գնդի կետերի օգնությամբ:

Հետևյալ թեորեմը հանդիսանում է որոշակի կետի՝ Ֆերմա--Տորիչելիի խնդրի լուծումների բազմությանը պատկանելու չափանիշ:

\textbf{Թեորեմ 2.4} ([7], [10]): \textit{Դիցուք $x_0, x_1, \dots, x_n$-ը տարածության մեջ կետեր են և $x_0 \ne x_i$՝ $i = 1, \dots, n$-ի համար: Այդ դեպքում $x_0$-ն Ֆերմա--Տորիչելիի կետ է $A = \{x_1, \dots, x_n\}$-ի համար այն և միայն այն դեպքում, երբ յուրաքանչյուր $x_i - x_0$ վեկտորի համար ($i = 1, \dots, n$) գոյություն ունի այնպիսի $\varphi_i$ նորմավորող ֆունկցիոնալ, որ $\sum_{i=1}^n \varphi_i = 0$:}

Ենթադրենք, որ գտնվել է մի կետ, որը հանդիսանում է լուծում $A$ բազմության համար: Այժմ թեորեմ 2.4-ի ֆունկցիոնալների օգնությամբ կարելի է կառուցել ամբողջ $\operatorname{ft}(A)$ բազմությունը: Դրա համար ներմուծենք նոր օբյեկտ:

\textbf{Սահմանում 2.5.} Դիցուք տրված են $\varphi \in X^*$ ֆունկցիոնալը և $x \in X$ կետը: Սահմանենք կոն.
\begin{equation*}
    C(x, \varphi) = x - \{a : \varphi(a) = \|a\|\}:
\end{equation*}
Այժմ մեկնաբանենք 
$$A_\varphi = \{a : \varphi(a) = \|a\|\}$$
բազմության կառուցվածքը (երկրաչափական իմաստը)։
Նկատենք, որ $a = 0$-ի համար՝
$$\varphi(0) = \varphi(0 \cdot 0) = 0 \cdot \varphi(0) = 0 \text{ և } \|0\| = 0 \Rightarrow \varphi(0) = \|0\|,$$ հետևաբար $a \in A_\varphi$։
Դիցուք այժմ $a \in A_\varphi$ և $a \neq 0$։ Դիտարկենք $x = \frac{a}{\|a\|}$ միավոր վեկտորը ($\|x\| = 1$)։ Այդ դեպքում $a = \|a\|x$ և, հետևաբար՝
$$\|a\| = \varphi(a) = \varphi(\|a\| \cdot x) = \|a\| \varphi(x) \Rightarrow \varphi(x) = 1,$$
այլ խոսքով, եթե $a \in A_\varphi$ և $a \neq 0$, ապա $a$-ով ուղղված միավոր $x$ վեկտորն այնպիսին է, որ $\varphi(x) = 1$։ Հակառակը, եթե $x$ միավոր վեկտորն այնպիսին է, որ $\varphi(x) = 1$ ($\|x\| = 1$), ապա $\forall t > 0$ թվի համար $a = t \cdot x \in A_\varphi$։
$$\varphi(a) = \varphi(t x) = t \, \varphi(x) = t = |t| = |t| \|x\| = \|t \cdot x\| = \|a\|$$

Սա նշանակում է, որ այդ $x$ վեկտորով անցնող ամբողջ ճառագայթն ամբողջությամբ պատկանում է $A_\varphi$-ին, այսինքն $A_\varphi$-ն սկիզբ է առնում զրոյից և իրենից ներկայացնում է կոն, որի ճառագայթներն անցնում են միավոր գնդի այնպիսի կետերով, որոնում $\varphi$ ֆունկցիոնալը <<հասնում է>> նորմի սահմանին (շոշափում է միավոր գունդը)։

Եթե $A_\varphi$-ի կառուցման մեջ մասնակցող $\varphi$ գծային ֆունկցիոնալն այնպիսին է, որ $\|\varphi\| < 1$, ապա $\|x\| = 1$ դեպքում՝
$$\varphi(x) \le |\varphi(x)| \le \|\varphi\| \|x\| = \|\varphi\| < 1,$$
այսինքն $\varphi(x) = 1$ տեղի ունենալ չի կարող, ուստի այդպիսի $A_\varphi$-ն պարպարունակում է միայն զրոյական կետ։

Դիցուք $\|\varphi\| = 1$։ Այն այն սահմանային դեպքն է, երբ $\varphi(x) = 1$ հավասարումը տեղի է ունենում միավոր սֆերայի միայն այն կետերում, որոնք հանդիսանում են հենման (շոշափման) կետեր։

Եթե $\|\varphi\| > 1$, ապա $\varphi(x) = 1$ գիծը (հարթությունը) կտրում է միավոր գունդը (անցնում է նրա ներսով), հետևաբար, այս դեպքում $A_\varphi$ բազմությունը ներկայացնում է միավոր սֆերայից և գնդից դուրս ձգվող կոն, որի հատույթը միավոր սֆերայի հետ տալիս է $\varphi(x) = 1$ հարթության և սֆերայի հատման գիծը։


\textbf{Թեորեմ 2.6} ([7]): \textit{Դիցուք $A = \{x_1, \dots, x_n\}$-ը տարածության մեջ կետեր են և $p \in \operatorname{ft}(A) \setminus A$: Ըստ թեորեմ 2.4-ի՝ յուրաքանչյուր $x_i - p$ վեկտորի համար ($i = 1, \dots, n$) գոյություն ունի այնպիսի $\varphi_i$ նորմավորող ֆունկցիոնալ, որ $\sum_{i=1}^n \varphi_i = 0$: Այդ դեպքում $\operatorname{ft}(A) = \bigcap_{i=1}^n C(x_i, \varphi_i)$:}

2.4 և 2.6 թեորեմները կազմում են Ֆերմա--Տորիչելիի խնդրի լուծման որոնման երկրաչափական մեթոդը: Այս թեորեմի կիրառությունը նկարագրելու համար ավելի մանրամասն դիտարկենք վեցանկյուն նորմում հավասարակողմ եռանկյան գագաթներով օրինակը (նկար 2):

Նախ և առաջ, գտնենք գոնե մեկ լուծում: Վերցնենք $p = \frac{1}{3}(x_1 + x_2 + x_3)$ և թեորեմ 2.4-ի օգնությամբ ապացուցենք, որ $p \in \operatorname{ft}(A)$: Դրա համար դիտարկենք $x_1 - p, x_2 - p, x_3 - p$ վեկտորները և կառուցենք դրանց նորմավորող ֆունկցիոնալները:Հենող ֆունկցիոնալ ասելով հասկանում ենք այնպիսի գծային ֆունկցիոնալ, որը երկրաչափորեն «հպվում» (շոշափում) է միավոր գնդին նրա սահմանային կետում կամ նիստում և ապահովում է, որ ամբողջ գունդը մնա այդ ֆունկցիոնալով սահմանափակված հարթության մի կողմում։ Այդ հենող ֆունկցիոնալը պետք է բավարարի հետևյալ պայմաններին.
1) Նորմը հավասար է $1'$ $\|\varphi\| = 1$։

2) Միավոր գնդի ցանկացած $x$ կետի համար՝$$\varphi(x) \leq 1$$ 
3)Միավոր սֆերայի վրա գոյություն ունի գոնե մեկ այնպիսի $x_0$ կետ ($\|x_0\|=1$), որտեղ $\varphi(x)$ ֆունկցիոնալի արժեքը հավասար է իր բացարձակ արժեքի առավելագույնին՝
$$\varphi(x_0) = 1$$
Քանի որ դրանք ընկած են հարթեցումների ներքին կետերի հետ միևնույն ուղղությունների վրա, ապա $x_i - p$ վեկտորներից յուրաքանչյուրի համար գոյություն ունի ճիշտ մեկ $\varphi_i$ ֆունկցիոնալ, որի մակարդակի գիծը պարունակում է համապատասխան հարթեցումը:Վեցանկյան $i$-դ կողմը որոշող ֆունկցիոնալը նշանակենք $\tilde{\varphi}_i$-ով: Այդ դեպքում կողմը կորոշվի $\tilde{\varphi}_i(x) = 1$ առնչությամբ: Միավոր գնդի ցանկացած $t$ վեկտորի համար կունենանք $\tilde{\varphi}_i(t) \leq 1$: Քանի որ $\tilde{\varphi}_i(t) \leq \|\tilde{\varphi}_i\|$ և ինչ-որ $\|x\| = 1$-ի համար $\tilde{\varphi}_i(x) = 1$, ապա $\|\tilde{\varphi}_i\| = 1$: Դա նշանակում է, որ $\tilde{\varphi}_i$-ն հենող ֆունկցիոնալ է:

Մյուս կողմից, եթե $\varphi_i$-ն $x_i-p$ վեկտորի նորմավորող ֆունկցիոնալ է, ապա $\|\varphi_i\| = 1$ և $\varphi_i(x_i-p) = \|x_i-p\|$: Այդ դեպքում $x_0^{(i)} = \frac{x_i-p}{\|x_i-p\|}$ տարրի համար կունենանք $\|x_0^{(i)}\| = 1$ և $\varphi_i(x_0^{(i)}) = 1$:

Պարզ է, որ $\|t\| \leq 1$ կետերի (վեկտորների) համար՝
$$\varphi_i(t) \leq |\varphi_i(t)| \leq \|\varphi_i\| = 1,$$
այսինքն՝ գնդին պատկանող կետերի համար՝
$$\varphi_i(t) \leq 1$$

Նկատենք, որ $\varphi_i$-ն չի մտնում գնդից ներս: Ենթադրենք դա այդպես չէ, այսինքն գոյություն ունի $y_0$ տարր (վեկտոր) այնպես, որ $\|y_0\| < 1$, բայց $\varphi_i(y_0) = 1$: Վերցնենք $t > 1$ թիվն այնպես, որ $z = t y_0$- համար $\|t y_0\| < 1$: Այդ դեպքում $\varphi_i(t y_0) = t \, \varphi_i(y_0) = t > 1$: Ստացվեց, որ գնդին պատկանող $z = t y_0$ տարրի (վեկտորի) համար $\varphi_i(z) > 1$: Ստացվեց հակասություն: Այսպիսով, $\varphi_i(x) = 1$ ուղիղն անցնում է $x_0^{(i)}$ կետով և չի մտնում միավոր վեցանկյան մեջ: Քանի որ $\tilde{\varphi}_i(x) = 1$ ուղիղն ևս անցնում է $x_0^{(i)}$ կետով, ուստի $\varphi_i = \tilde{\varphi}_i$: Հարթեցումները հավասարահեռ են կոորդինատների սկզբնակետից և որոշում են հավասարակողմ եռանկյուն, հետևաբար, դրանցով կառուցված ֆունկցիոնալների գումարը հավասար է զրոյի, և թեորեմի պայմանը կատարված է: Քանի որ մենք աշխատում ենք երկչափ հարթության վրա, ցանկացած $v$ վեկտոր ունի կոորդինատներ՝ $v = (x, y)$։ Հարթության վրա գծային ֆունկցիոնալը այլ բան չէ, քան երկու թվերի $(a, b)$ զույգ, որը $v = (x, y)$ վեկտորը բազմապատկում է հետևյալ կերպ.
$$\varphi(v) = a \cdot x + b \cdot y$$

Ուստի մեր երեք ֆունկցիոնալները որոշվում են իրենց գործակիցներով՝
$$\varphi_1(x, y) = a_1 x + b_1 y, \quad \varphi_2(x, y) = a_2 x + b_2 y$$
$$\varphi_3(x, y) = a_3 x + b_3 y$$

Այդ դեպքում $\varphi_1 + \varphi_2 + \varphi_3 = 0$, ինչը նշանակում է, որ՝
$$\begin{cases} 
a_1 + a_2 + a_3 = 0 \\ 
b_1 + b_2 + b_3 = 0 
\end{cases}$$
Եթե $120^{\circ}$-ով պտտում ենք առաջին ֆունկցիոնալը` $\varphi_1$-ը, ստանում ենք երկրորդ ֆունկցիոնալի ($\varphi_2$) գործակիցները` $(a_2; b_2)$, իսկ $240^{\circ}$-ով պտտելիս` երրորդը ($\varphi_3$, այսինքն` $(a_3; b_3)$): Այսինքն $(a_i, b_i)$ զույգերը ձևավորում են ճշգրիտ նույն երկարությամբ $3$ վեկտորներ, որոնք միմյանց նկատմամբ շեղված են $120^{\circ}$-ով:

Այդ դեպքում`
\begin{equation*}
(a_1, b_1) = (r \cos \theta, r \sin \theta)
\end{equation*}
\begin{equation*}
(a_2, b_2) = (r \cos(\theta + 120^{\circ}), r \sin(\theta + 120^{\circ}))
\end{equation*}
\begin{equation*}
(a_3, b_3) = (r \cos(\theta + 240^{\circ}), r \sin(\theta + 240^{\circ}))
\end{equation*}

Այդ դեպքում դժվար չէ համոզվել, որ $a_1 + a_2 + a_3 = 0$ և $b_1 + b_2 + b_3 = 0$:

\begin{figure}[htbp]
\centering
\includegraphics[width=0.75\textwidth]{images/2.png}
\caption{Լուծումների լրիվ բազմության կառուցումը $x_1, x_2, x_3$ կետերի համար վեցանկյուն նորմով հարթության վրա}
\label{fig:2}
\end{figure}

Այժմ օգտվենք թեորեմ 2.6-ից, որպեսզի գտնենք բոլոր լուծումները: Կառուցենք $\varphi_2$ ֆունկցիոնալով և $x_2 - p$ վեկտորով տրվող կոնը: Քանի որ ֆունկցիոնալը հանդիսանում է նորմավորող հարթեցման բոլոր կետերի համար, $\{a : \varphi_2(a) = \|a\|\}$ բազմությունն իրենից ներկայացնում է կոորդինատների սկզբնակետից ելնող և հարթեցման կետերով անցնող ճառագայթների հավաքածու, այսինքն՝ անկյուն, որի կողմերը պարունակում են երկու հարևան գագաթներ: Այժմ անկյունն անդրադարձնենք կոորդինատների սկզբնակետի նկատմամբ և կատարենք զուգահեռ տեղափոխություն այնպես, որ անկյան գագաթը հայտնվի $x_2$ կետում: Մյուս երկու կոնների նմանատիպ կառուցումից հետո կստանանք, որ դրանց հատումն ամբողջ $x_1 x_2 x_3$ եռանկյունն է:

\textbf{Դիտողություն 2.7.} Թեորեմ 2.6-ի պնդումը կախված չէ $p$ կետի և $\varphi_i$ ֆունկցիոնալների ընտրությունից:

Երկրաչափական մեթոդը տալիս է Ֆերմա--Տորիչելիի խնդրի լուծումների ընդհանուր նկարագրությունը տարբեր տրված բազմությունների համար. դա որոշակի կոների հատումն է՝ այդ բազմության կետերում տեղակայված գագաթներով: Այս դիտարկումը թույլ է տալիս ձևակերպել հետևյալ պնդումը.

\textbf{Պնդում 2.8} ([13]): \textit{Դիցուք տարածության մեջ $A = \{x_1, \dots, x_{2k+1}\}$ բազմության կետերը դասավորված են մեկ ուղղի վրա՝ իրենց համարակալման կարգով: Այդ դեպքում $\operatorname{ft}(A) = \{x_{k+1}\}$: Եթե $A = \{x_1, \dots, x_{2k}\}$, ապա $\operatorname{ft}(A) = x_k x_{k+1}$:}

\section{Միակության չափանիշը}

\subsection{Չափանիշ \texorpdfstring{$d$}{d} չափողականությամբ տարածությունում \texorpdfstring{$n$}{n} կետերի համար}

Դիցուք $X$-ը $d$ չափողականությամբ Մինկովսկու տարածություն է: $X$-ում դնենք Ֆերմա--Տորիչելիի խնդիրը $n$ կետերի համար: Եթե $x_i \in X$, $1 \le i \le n$, ապա լուծում է կոչվում $\operatorname{ft}(x_1, \dots, x_n)$ բազմությունը, որի մեջ մտնում են բոլոր այն $x \in X$ կետերը, որոնց վրա հասնում է $\sum_{i=1}^n |x x_i|$ ֆունկցիայի մինիմումը:

Դիտարկենք $X$ տարածության $S$ միավոր սֆերան:

\textbf{Սահմանում 3.1.} Միավոր սֆերայի \textit{նիստ} կոչվում է նրա հատումը որևէ հենման հիպերհարթության հետ: Եթե նիստի կետերի գծային թաղանթը $X$-ի $k$ չափողականությամբ ենթատարածություն է, ապա այդպիսի նիստը կոչվում է $(k-1)$-չափանի:

\textbf{Սահմանում 3.2.} Վերցնենք Մինկովսկու տարածության $S$ միավոր սֆերայի նիստերի վերջավոր բազմություն: Դրանցից յուրաքանչյուրի համար ընտրենք $\pi_i$ հենման հիպերհարթություն, որը հատվում է $S$-ի հետ միայն այդ նիստով: Դիցուք $\pi_i$ հիպերհարթությունը որոշում է որոշակի $\varphi_i$ գծային ֆունկցիոնալի $\varphi_i = 1$ մակարդակի մակերևույթը: Եթե գտնվի հենման հիպերհարթությունների այնպիսի հավաքածու, որ կառուցված ֆունկցիոնալների գումարը հավասար լինի զրոյի, ապա նիստերի այդպիսի բազմությունը կանվանենք \textit{համաձայնեցված}:

Սահմանումից հետևում է, որ եթե համաձայնեցված բազմության յուրաքանչյուր նիստից վերցնենք մեկական կետ, ապա կստանանք կետերի հավաքածու, որի համար $x = 0$-ն հանդիսանում է Ֆերմա--Տորիչելիի խնդրի լուծումներից մեկը: Կախված բազմության մեջ մտնող նիստերի չափայնությունից և դրանց փոխադարձ դասավորությունից՝ կարելի է ստանալ լուծումների տարբեր տիպեր, այդ թվում՝ միակ և ոչ միակ: Պարզվում է, որ $X$ տարածությունում ոչ միակ լուծումով կետերի բազմության որոնման խնդիրը համարժեք է այդ տարածությունում որոշակի հատկություններով նիստերի համաձայնեցված բազմության գոյությանը:

\textbf{Թեորեմ 3.3.} \textit{Դիցուք $X$-ը $d$ չափողականությամբ Մինկովսկու տարածություն է:}

\textit{Եթե $n \ge 3$ կենտ է, ապա $X$ տարածությունում կգտնվեն $n$ կետեր, որոնց համար Ֆերմա--Տորիչելիի խնդրի լուծումը միակը չէ այն և միայն այն դեպքում, երբ $X$-ում գոյություն ունի $S$ միավոր սֆերայի $n$ նիստերից բաղկացած համաձայնեցված բազմություն, որի համար կատարվում է հետևյալ պայմանը.}
\begin{itemize}[label=\textbullet]
    \item \textit{Այդ համաձայնեցված բազմության մեջ մտնող բոլոր նիստերի գծային թաղանթների հատումը պարունակում է ուղիղ:}
\end{itemize}

\textit{Եթե $n \ge 4$ զույգ է, ապա $X$ տարածությունում Ֆերմա--Տորիչելիի խնդրի լուծումը միակն է մեկ ուղղի վրա չգտնվող ցանկացած $n$ կետերի համար այն և միայն այն դեպքում, երբ նորմը խիստ ուռուցիկ է:}

\textbf{Ապացույց:} Դիցուք $n$-ը կենտ է:

\textbf{Բավարարություն:} Ապացուցենք, որ եթե այդպիսի $A$ համաձայնեցված բազմություն գոյություն ունի, ապա գոյություն ունի կետերի հավաքածու, որոնց համար Ֆերմա--Տորիչելիի խնդրի լուծումը ոչ միակ է: Յուրաքանչյուր նիստից վերցնենք մեկական ներքին կետ: Դիցուք դրանք $x_1, \dots, x_n$ կետերն են: Ըստ թեորեմ 2.4-ի՝ $x = 0$ կետը հանդիսանում է լուծումներից մեկը: Ըստ թեորեմ 2.6-ի՝ լրիվ լուծումն այն կոների հատումն է, որոնք ելնում են գագաթներից և պարունակում են $x = 0$ կետը: Ընդ որում $k$-չափանի նիստից ելնում է $(k+1)$-չափանի կոն: Ըստ պայմանի՝ գոյություն ունի $l$ ուղիղ, որը պատկանում է $A$-ի ցանկացած նիստի գծային թաղանթին, այսինքն՝ ցանկացած կոնի գծային թաղանթին: Այդ դեպքում այդ կոներից յուրաքանչյուրը պարունակում է $x = 0$ կետի որոշակի միաչափ շրջակայք $l$ ուղղի վրա: Հետևաբար, բոլոր կոների հատումը նույնպես պարունակում է զրոյի շրջակայք: Լուծումը ոչ միակն է:

\textbf{Անհրաժեշտություն:} Դիցուք $X$ տարածությունում գտնվել են $n$ կետեր՝ $x_1, \dots, x_n$, որոնց համար Ֆերմա--Տորիչելիի խնդրի լուծումը ոչ միակ է: Դիցուք $p$ կետը մտնում է լուծման մեջ: Այդ դեպքում $x_i - p$ կետերի համար լուծումը նույնպես ոչ միակ է, և $x = 0$-ն լուծումներից մեկն է: Ըստ թեորեմ 2.6-ի՝ $\operatorname{ft}(x_1 - p, \dots, x_n - p)$ բազմությունը $x_i - p$ կետերից ելնող որոշակի $C_1, \dots, C_n$ կոների հատումն է: Դիտարկենք այն նորմավորող ֆունկցիոնալները, որոնցով կառուցված են կոները, այն է՝ միավոր սֆերայի հատումները $\varphi_i = 1$ հենման հիպերհարթությունների հետ: Կստանանք $n$ նիստեր, որոնք ակնհայտորեն կազմում են համաձայնեցված բազմություն: Քանի որ լուծումը ոչ միակ է և հանդիսանում է կոների հատում, ապա այն պարունակում է որոշակի ոչ դատարկ հատված, որը պարունակում է $x = 0$-ն: Յուրաքանչյուր կոն նույնպես պարունակում է այդ հատվածը: Այդ դեպքում կոնի գծային թաղանթում ընկած է այդ հատվածը պարունակող ուղիղը: Հետևաբար, բոլոր նիստերի գծային թաղանթների հատումը պարունակում է ուղիղ:

Դիցուք $n$-ը զույգ է:

Դիցուք $n$-ը զույգ է և նորմը խիստ ուռուցիկ է: Այդ դեպքում Ֆերմա--Տորիչելիի խնդրի լուծումը ցանկացած $n$ կետերի համար այդ կետերում գագաթներով միաչափ կոների հատումն է: Քանի որ գագաթները չեն գտնվում մեկ ուղղի վրա, ապա կոները կարող են հատվել միայն մեկ կետով, այսինքն՝ լուծումը միշտ միակ է:

Դիցուք միավոր սֆերան պարունակում է նիստ, որը կետ չէ: Նրանում վերցնենք ցանկացած $\frac{n}{2}$ տարբեր ներքին կետեր՝ $x_1, \dots, x_{n/2}$: Դիտարկենք լուծումների բազմությունը $\pm x_1, \dots, \pm x_{n/2}$ հավաքածուի համար: Կետերից յուրաքանչյուրի համար տանք ֆունկցիոնալ, որի մակարդակի մակերևույթը հանդիսանում է այն նիստը որոշող հենման հիպերհարթությունը, որին պատկանում է կետը: Քանի որ նիստերը հակադիր են և պարունակում են կետերի միևնույն քանակ, ապա ֆունկցիոնալների գումարը հավասար է զրոյի: Ըստ թեորեմ 2.4-ի՝ $x = 0$-ն լուծումներից մեկն է: Լրիվ լուծումը հանդիսանում է կոների հատում: Բոլոր կոներն ունեն միևնույն չափողականությունից, ընկած են նույն չափողականության մեկ ենթատարածությունում և պարունակում են $x = 0$ որպես ներքին կետը: Այսպիսով, բոլոր կոների հատմանը պատկանում է համապատասխան չափողականությամբ զրոյի որոշակի շրջակայք: Լուծումը ոչ միակ է:

Թեորեմն ապացուցված է:

\textbf{Դիտողություն 3.4.} Քանի որ առավելագույն չափողականությամբ նիստի գծային թաղանթն ամբողջ տարածությունն է, ապա թեորեմի պայմանում կարելի է դիտարկել միայն ոչ առավելագույն նիստերի հատումը: Այսպիսով, եթե համաձայնեցված բազմությունը բաղկացած է միմիայն առավելագույն նիստերից, ապա լրացուցիչ պայմանը կատարվում է ավտոմատ կերպով:

Եթե տարածության $d$ չափողականությունը հավասար է 2-ի կամ 3-ի, ապա համարժեք պայմանը հաջողվում է ճշգրտել: Առավել մանրամասն դիտարկենք այս երկու դեպքերը:

\subsection{Հարթության վրա \texorpdfstring{$n$}{n} կետերի համար որոշված չափանիշ}

Դիցուք $X$-ը Մինկովսկու հարթություն է: $X$-ում դրված է Ֆերմա--Տորիչելիի խնդիրը $n$ կետերի համար: $S$ միավոր շրջանագիծն ունի միայն երկու տիպի նիստեր: Դրանք կանվանենք կետեր և հարթեցումներ:

\textbf{Պնդում 3.5.} \textit{Եթե $X$-ում գոյություն ունի միավոր շրջանագծի նիստերի համաձայնեցված բազմություն՝ բաղկացած $n$ հարթեցումներից, ապա այդ հարթության մեջ կգտնվեն $x_1, \dots, x_n$ կետեր, որոնց համար $\operatorname{ft}(x_1, \dots, x_n)$ բազմությունը չվերասերված բազմանկյուն է:}

\textbf{Ապացույց:} Դիցուք $x_1, \dots, x_n$-ը համաձայնեցված բազմության հարթեցումների ներքին կետեր են: Այդ դեպքում, վերցնելով $\varphi_1, \dots, \varphi_n$ ֆունկցիոնալները սահմանում 3.2-ից, ըստ թեորեմ 2.4-ի ստանում ենք, որ $p = 0$-ն պատկանում է $\operatorname{ft}(x_1, \dots, x_n)$-ին: Ըստ թեորեմ 2.6-ի՝ $\operatorname{ft}(x_1, \dots, x_n)$-ը $x_1, \dots, x_n$ կետերից ելնող կոների հատումն է: Քանի որ բոլոր ֆունկցիոնալները պարունակում են հարթեցումներ, բոլոր կոները չվերասերված անկյուններ են: Ընդ որում դրանցից յուրաքանչյուրը պարունակում է զրոյի որոշակի լիաչափ շրջակայք այն պատճառով, որ $x_i$ կետերը հարթեցումների սահմաններ չեն: Հետևաբար, կոները հատման մեջ կազմում են բազմանկյուն:

\begin{figure}[htbp]
\centering
\includegraphics[width=0.65\textwidth]{images/3.png}
\caption{Ոչ միակության առաջին պայմանը $n = 3$ դեպքում}
\label{fig:3}
\end{figure}

\textbf{Պնդում 3.6.} \textit{Եթե $X$-ում գոյություն ունի միավոր շրջանագծի նիստերի համաձայնեցված բազմություն՝ բաղկացած $n-1$ հարթեցումներից և մեկ կետից, ապա այդ հարթության մեջ կգտնվեն $x_1, \dots, x_n$, որոնց համար $\operatorname{ft}(x_1, \dots, x_n)$ լուծումների բազմությունը չվերասերված հատված է:}

\textbf{Ապացույց:} Դիցուք $x_1, \dots, x_{n-1}$ կետերը հարթեցումների ներքին կետեր են, իսկ $x_n$-ը՝ կետ համաձայնեցված բազմությունից: Նախորդ ապացուցման նմանությամբ ստանում ենք, որ $p = 0$-ն պատկանում է $\operatorname{ft}(x_1, \dots, x_n)$-ին: $x_1, \dots, x_{n-1}$-ից ելնող կոները պարունակում են զրոյի շրջակայք, իսկ $x_n$ կետից ելնում է կոորդինատների սկզբնակետով անցնող ճառագայթ: Հետևաբար, հատման մեջ ստանում ենք չվերասերված հատված:

\textbf{Պնդում 3.7.} \textit{Դիցուք $2 \le k \le n - 1$: Եթե $X$-ում գտնվի միավոր շրջանագծի նիստերի համաձայնեցված բազմություն՝ բաղկացած $k$ հարթեցումներից և $n - k$ կետերից, ընդ որում կետերի գծային թաղանթն ուղիղ է, ապա այդ հարթության մեջ կգտնվեն $x_1, \dots, x_n$,կետեր  որոնց համար $\operatorname{ft}(x_1, \dots, x_n)$ լուծումների բազմությունը չվերասերված հատված է:}

\textbf{Ապացույց:} Դիցուք $x_1, \dots, x_k$-ը հարթեցումների ներքին կետեր են, իսկ $x_{k+1}, \dots, x_n$-ը՝ կետեր համաձայնեցված բազմությունից: Սահմանում 3.2-ի $\varphi_1, \dots, \varphi_n$ ֆունկցիոնալները բավարարում են թեորեմ 2.4-ին, և $x = 0$ կետը պատկանում է $\operatorname{ft}(x_1, \dots, x_n)$-ին: Ըստ թեորեմ 2.6-ի՝ $\operatorname{ft}(x_1, x_2, x_3)$ լուծումը հանդիսանում է կոորդինատների սկզբնակետը պարունակող և միևնույն ուղղի վրա ընկած $k$ ճառագայթների, ինչպես նաև $n - k$ չվերասերված անկյունների հատում, որոնցից յուրաքանչյուրը պարունակում է զրոյի շրջակայք: Ստանում ենք հատված: 

\begin{figure}[htbp]
\centering
\includegraphics[width=0.65\textwidth]{images/4.png}
\caption{Ոչ միակության երրորդ պայմանը $n = 3$ դեպքում}
\label{fig:4}
\end{figure}

\textbf{Թեորեմ 3.8.} \textit{Եթե $n \ge 3$ կենտ է, ապա նորմավորված հարթությունում Ֆերմա--Տորիչելիի խնդրի լուծումը միակ է ցանկացած $n$ կետերի համար այն և միայն այն դեպքում, երբ միավոր շրջանագծի նիստերի ցանկացած համաձայնեցված բազմություն բաղկացած է ոչ ավելի, քան $n - 2$ հարթեցումներից, և նիստ-կետերի գծային թաղանթի չափողականությունը հավասար է 2-ի:}

\textbf{Ապացույց:} Երկու պնդումներում էլ անհրաժեշտությունը բխում է 3.5, 3.6 և 3.7 պնդումներից: Ապացուցենք բավարարությունը: Դիցուք $n$-ը կենտ է և գտնվել են $x_1, \dots, x_n$ կետեր, որոնց համար Ֆերմա--Տորիչելիի խնդրի լուծումների բազմությունը ոչ միակ է: Քանի որ լուծումը կոների հատում է, ապա դա կա՛մ բազմանկյուն է, կա՛մ հատված: Բազմանկյան դեպքում բոլոր կոները չվերասերված անկյուններ են: Դիտարկենք կոները որոշող ֆունկցիոնալները: Դրանց մակարդակի գծերը պարունակում են միավոր շրջանագծի հարթեցումներ, որոնք ակնհայտորեն կազմում են նիստերի համաձայնեցված բազմություն:

Հատվածը կարող է ստացվել գոնե մեկ չվերասերված անկյան և մեկ ուղղի վրա ընկած ճառագայթների հատմամբ: Եթե բոլոր կոները ճառագայթներ են, ապա բոլոր $x_i$ կետերն ընկած են մեկ ուղղի վրա, բայց ըստ պնդում 2.8-ի՝ այդ դեպքում լուծումը միակ է: Նմանապես դիտարկենք կոները որոշող ֆունկցիոնալները: Քանի որ բոլոր չվերասերված կոներն ընկած են մեկ ուղղի վրա, ապա միավոր շրջանագծի և ֆունկցիոնալների մակարդակի գծերի հատման կետերը նույնպես ընկած են մեկ ուղղի վրա: Այսինքն՝ համաձայնեցված բազմության կետերի գծային թաղանթը հավասար է 1-ի: 

\subsection{Եռաչափ տարածությունում \texorpdfstring{$n$}{n} կետերի համար որոշված չափանիշ}

Դիցուք $X$-ը Մինկովսկու եռաչափ տարածություն է: $X$-ում դրված է Ֆերմա--Տորիչելիի խնդիրը $n$ կետերի համար:

\textbf{Թեորեմ 3.9.} \textit{Եթե $n \ge 3$ կենտ է, ապա Մինկովսկու $X$ եռաչափ տարածությունում կգտնվեն $n$ կետեր, որոնց համար Ֆերմա--Տորիչելիի խնդրի լուծումը ոչ միակ է, այն և միայն այն դեպքում, երբ $X$-ում գոյություն ունի $S$ միավոր սֆերայի $n$ նիստերից բաղկացած համաձայնեցված բազմություն, որի համար կատարվում են հետևյալ պայմանները.}

\begin{itemize}[label=\textbullet]
    \item \textit{Այդ համաձայնեցված բազմության մեջ մտնող նիստ-կետերի գծային թաղանթների հատումը պարունակում է ուղիղ,}
    \item \textit{Այդ համաձայնեցված բազմության մեջ մտնող նիստ-կետի և նիստ-հատվածի ցանկացած զույգի համար դրանց գծային թաղանթի չափողականությունը հավասար է 2-ի,}
    \item \textit{Այդ համաձայնեցված բազմության մեջ մտնող նիստ-հատվածների գծային թաղանթների հատումը պարունակում է ուղիղ:}
\end{itemize}

\textbf{Ապացույց:} Դիցուք $n$-ը կենտ է:

\textbf{Բավարարություն:} Դիցուք գտնվել է $n$ նիստերից բաղկացած համաձայնեցված բազմություն, որի համար բոլոր պայմանները կատարվում են: Կառուցենք Ֆերմա--Տորիչելիի խնդրի լուծումը տվյալ նիստերի ներքին կետերի հավաքածուի համար: Ըստ թեորեմ 2.4-ի՝ $x = 0$ կետը մտնում է լուծման մեջ: Լրիվ լուծումն այն կոների հատումն է, որոնք ելնում են գագաթներից և պարունակում են $x = 0$-ն որպես ներքին կետ:

Դիցուք համաձայնեցված բազմությունն ունի նիստ-կետեր: Այդպիսի նիստի գծային թաղանթն այն ուղիղն է, որը պարունակում է հենց այդ կետը և $x = 0$-ն: Ելնելով համաձայնեցված բազմության վրա դրվող առաջին պայմանից՝ նրա բոլոր կետերի համար այդ ուղիղը համընկնում է: Դիցուք դա $l$ ուղիղն է: Նիստ-կետերից ելնող բոլոր միաչափ կոների հատումը պարունակում է ոչ դատարկ հատված, որն ընկած է $l$-ում և պարունակում է $x = 0$-ն որպես ներքին կետ:

Երկրորդ տիպի նիստերի գծային թաղանթներն այն հարթություններն են, որոնցում ընկած են տվյալ տարրերից ելնող երկչափ կոները: Ըստ պայմանի՝ յուրաքանչյուր այդպիսի հարթություն պարունակում է բոլոր նիստ-կետերը, իսկ հետևաբար նաև $l$ ուղիղը, որի վրա դրանք ընկած են: Գոյություն ունի $l$-ում ընկած և $x = 0$-ն որպես ներքին կետ պարունակող ոչ դատարկ հատված, որն ընկած է երկչափ կոներից յուրաքանչյուրում: Ստանում ենք, որ բոլոր միաչափ և երկչափ կոների հատումը պարունակում է առնվազն հատված: Եթե համաձայնեցված բազմության մեջ չկան նիստ-կետեր, ապա ըստ պայմանի՝ որպես $l$ ուղիղ վերցնենք հարթությունների հատման մեջ ընկած ցանկացած ուղիղ:

Քանի որ եռաչափ կոների հատման մեջ մտնում է $x = 0$ կետի որոշակի լիաչափ շրջակայք, ապա դրանց հատումն ստացված հատվածի հետ ոչ միակ է: Այսինքն՝ բոլոր կոների հատումը պարունակում է առնվազն հատված: Բավարարությունն ապացուցված է:

\textbf{Անհրաժեշտություն:} Դիցուք $X$ տարածությունում գտնվել են $n$ կետեր՝ $x_1, \dots, x_n$, որոնց համար Ֆերմա--Տորիչելիի խնդրի լուծումը ոչ միակ է: Դիցուք $p$ կետը մտնում է լուծման մեջ: Այդ դեպքում $x_i - p$ կետերի համար լուծումը նույնպես ոչ միակ է, և $x = 0$-ն լուծումներից մեկն է: Ըստ թեորեմ 2.6-ի՝ $\operatorname{ft}(x_1 - p, \dots, x_n - p)$ բազմությունը $x_i - p$ կետերից ելնող որոշակի $C_1, \dots, C_n$ կոների հատումն է: Դիտարկենք այն նորմավորող ֆունկցիոնալները, որոնցով կառուցված են կոները, այն է՝ միավոր սֆերայի հատումները $\varphi_i = 1$ հենման հարթությունների հետ: Կստանանք միավոր սֆերայի նիստերի համաձայնեցված բազմություն: Որոշենք, թե ինչ պայմաններ են դրվում տվյալ բազմության վրա:

Նիստ-կետերից ելնող բոլոր միաչափ կոները պետք է ընկած լինեն մեկ ուղղի վրա, հակառակ դեպքում դրանց հատումը կպարունակի միայն $x = 0$ կետը: Այլ խոսքով՝ այդ նիստերի գծային թաղանթների հատումը հենց այդ ուղիղն է: Դիցուք դա $l$ ուղիղն է:

Նիստ-հատվածներից ելնում են երկչափ կոներ: Քանի որ լուծումը ոչ միակ է, ուստի դրանց հատման մեջ ընկած է $x = 0$-ն պարունակող ոչ դատարկ հատված: Դիցուք նաև այդ հատվածը պատկանում է $l$ ուղղին, եթե համաձայնեցված բազմությունը պարունակում է կետեր և այդպիսի ուղիղ կառուցված է:

Այդ դեպքում ցանկացած նիստ-հատվածի գծային թաղանթը, այսինքն՝ համապատասխան կոնը պարունակող հարթությունը, պարունակում է այդ հատվածը պարունակող ուղիղը: Ստացանք, որ նիստ-հատվածների գծային թաղանթների հատումը պարունակում է ուղիղ: Նիստ-կետերի առկայության դեպքում դա կլինի $l$ ուղիղը: Ուրեմն, նիստ-հատվածը պարունակող հարթությունը պարունակում է նաև ցանկացած նիստ-կետ: Այսինքն՝ նիստ-կետի և նիստ-հատվածի ցանկացած զույգի գծային թաղանթն ունի 2 չափող: Համաձայնեցված բազմության վրա դրվող բոլոր պայմանները կատարված են: 

\section{Չափանիշի կիրառումը}

\subsection{\texorpdfstring{$\lambda$}{lambda}-նորմավորված հարթություններ}

Գործնականում 3.3 չափանիշը կիրառելը հեշտ չէ, քանի որ տրված նորմի համար պատասխանը տրվում է միավոր շրջանագծի հարթեցումների ընտրության ճանապարհով: Եթե այն դիտարկենք որպես ալգորիթմ, ապա դրա բարդությունն ուղղակիորեն կախված է դրանց քանակից, իսկ անվերջ քանակով հարթեցումներ ունեցող նորմերի առկայության պատճառով այն նույնիսկ հաշվարկելի չէ վերջավոր ժամանակում: Ցույց տանք դրա կիրառությունը ոչ խիստ ուռուցիկ նորմերի ամենատարածված ընտանիքի՝ $\lambda$-հարթությունների վրա:

\textbf{Սահմանում 4.1.} \textit{$\lambda$-նորմավորված հարթություն} է կոչվում այն հարթությունը, որի վրա նորմը տրված է կանոնավոր $2\lambda$-անկյունով:

\textbf{Պնդում 4.2.} \textit{$\lambda$-նորմավորված հարթության համար, երբ $\lambda \not\equiv 0 \pmod 3$, չի կատարվում ոչ միակության առաջին պայմանը:}

\textbf{Ապացույց:} Ենթադրենք, որ պայմանը կատարվում է, և կգտնվեն այնպիսի $\varphi_1, \varphi_2, \varphi_3$ ֆունկցիոնալներ, որոնց գումարը հավասար է զրոյի, և $\varphi_i = 1$ ուղիղները պարունակում են միավոր շրջանագծի հարթեցումներ: Դիցուք $l_i$-երը $\varphi_i = 1$ մակարդակի գծերը որոշող ուղիղներն են: Այդ ուղիղները որոշող վեկտորների գումարը հավասար է զրոյի, ինչպես նաև $d(0, l_1) = d(0, l_2) = d(0, l_3)$, քանի որ $l_i$-երը պարունակում են $2\lambda$-անկյան հարթեցումներ: Հետևաբար, մակարդակի գծերը կազմում են կանոնավոր եռանկյուն, ինչն անհնար է, երբ $\lambda \not\equiv 0 \pmod 3$: 

\textbf{Պնդում 4.3.} \textit{$\lambda$-նորմավորված հարթության համար, երբ $\lambda \not\equiv 0 \pmod 3$, չի կատարվում ոչ միակության երկրորդ պայմանը:}

\textbf{Ապացույց:} Ենթադրենք, որ պայմանը կատարվում է, և կգտնվեն այնպիսի $\varphi_1, \varphi_2, \varphi_3$ ֆունկցիոնալներ, որոնց գումարը հավասար է զրոյի, $\varphi_1 = 1$ և $\varphi_2 = 1$ ուղիղները պարունակում են միավոր շրջանագծի հարթեցումներ, իսկ $\varphi_3 = 1$ ուղիղը՝ գագաթ:

Դիցուք $l_i$-երը $\varphi_i = 1$ մակարդակի գծերը որոշող ուղիղներն են: Հարթությունը պտտենք այնպես, որ $l_3$ ուղիղը լինի զուգահեռ $x$ կոորդինատային առանցքին: Դիցուք $l_1$-ը և $l_2$-ը տրվում են $\alpha_1 x + \beta_1 y = 1$ և $\alpha_2 x + \beta_2 y = 1$ հավասարումներով, իսկ $l_3$-ը՝ $\gamma y = 1$ հավասարմամբ: Քանի որ $\sum_{i=1}^3 \varphi_i = 0$, ապա $\alpha_1 + \alpha_2 = 0$ և $\beta_1 + \beta_2 + \gamma = 0$: Այն բանից, որ $l_1$ և $l_2$ ուղիղներին պատկանում են բազմանկյան հարթեցումներ, ստանում ենք.
\begin{equation*}
    d(0, l_1) = d(0, l_2) \Rightarrow \alpha_1^2 + \beta_1^2 = \alpha_2^2 + \beta_2^2:
\end{equation*}
Հաշվի առնելով, որ $l_1$-ը և $l_2$-ը զուգահեռ չեն՝
\begin{equation*}
    -\alpha_1 = \alpha_2 = \alpha, \quad \beta_1 = \beta_2 = \beta, \quad \gamma = -2\beta:
\end{equation*}

Այսպիսով, $l_i$ ուղիղների հավասարումներն ընդունեցին հետևյալ տեսքը.
\begin{equation*}
    -\alpha x + \beta y = 1, \quad \alpha x + \beta y = 1, \quad -2\beta y = 1:
\end{equation*}
$l_i$ ուղիղներով կազմված եռանկյան մեջ $l_1$ և $l_2$ ուղիղների հատման կետից տանենք կիսորդ: Կիսորդն ընկած է օրդինատների առանցքի վրա: Հնարավոր է երկու դեպք. կա՛մ այն անցնում է $2\lambda$-անկյան հանդիպակաց կողմերի միջնակետերով, կա՛մ հակադիր գագաթների զույգով: Առաջին դեպքը հակասում է նրան, որ $l_3$ հենման ուղիղը չի պարունակում հարթեցում: Այդ դեպքում կիսորդն անցնում է այն գագաթով, որին շոշափում է $l_3$-ը: Քանի որ $l_3$-ը զուգահեռ է աբսցիսների առանցքին, ստանում ենք $2\lambda$-անկյան կողմերի և $l_3$ ուղղի կազմած անկյունների հավասարությունը:

$\alpha$ և $\beta$ գործակիցների վրա դնենք ևս մեկ պայման: Դիցուք $n$-ը բազմանկյան գագաթների քանակն է: Ունենք.
\begin{equation*}
    d(0, l_1) = d(0, l_2) = \cos\frac{\pi}{n} d(0, l_3) \Rightarrow \alpha^2 + \beta^2 = \frac{4}{\cos^2\frac{\pi}{n}} \beta^2:
\end{equation*}
Այստեղից՝
\begin{equation*}
    \frac{\alpha}{\beta} = \sqrt{\frac{4}{\cos^2\frac{\pi}{n}} - 1} = \sqrt{4\operatorname{tg}^2\frac{\pi}{n} + 3}:
\end{equation*}
Եթե $\xi_{12}, \xi_{13}, \xi_{23}$-ը $l_1, l_2, l_3$ ուղիղների միջև անկյուններն են, ապա $\xi_{13} = \xi_{23}$ և $\xi_{12} + \xi_{13} + \xi_{23} = \pi$: Քանի որ $\tan\xi_{13} = \tan\xi_{23} = \frac{\alpha}{\beta}$, ապա
\begin{equation*}
    \xi_{12} = \pi - 2\operatorname{arctg}\operatorname{tg}\xi_{13} = \pi - 2\operatorname{arctg}\sqrt{4\operatorname{tg}^2\frac{\pi}{n} + 3}:
\end{equation*}
Մյուս կողմից, $\xi_{12} = \pi - \frac{2\pi k}{n}$ որոշակի $k$-ի համար, $k \in \mathbb{N}, \ 1 < k < \frac{n}{2}$: Այսինքն՝
\begin{equation*}
    \frac{\pi k}{n} = \operatorname{arctg}\sqrt{4\operatorname{tg}^2\frac{\pi}{n} + 3}:
\end{equation*}
Ապացուցենք, որ այս հավասարությունը տեղի չի ունենում $k$-ի և $n$-ի այնպիսի բնական արժեքների դեպքում, որ $n \equiv \pm 1 \pmod 3$, և կհանգենք հակասության նախնական ենթադրության հետ: Դիցուք
\begin{equation*}
    K(n) = \frac{n}{\pi} \operatorname{arctg}\sqrt{4\operatorname{tg}^2\frac{\pi}{n} + 3}
\end{equation*}
\begin{equation*}
    \frac{n}{3} - \frac{1}{3} < k(n) < \frac{n}{3} + \frac{1}{3}
\end{equation*}
անհավասարությունը տեղի ունի բոլոր $n > 2$-ի համար: Պնդումն ապացուցված է: 

\textbf{Պնդում 4.4.} \textit{$\lambda$-նորմավորված հարթության համար չի կատարվում ոչ միակության երրորդ պայմանը:}

\textbf{Ապացույց:} Ենթադրենք, որ պայմանը կատարվում է: Դիցուք հարթեցումներից մեկը պատկանում է որոշակի ֆունկցիոնալի $\varphi = 1$ մակարդակի գծին, ընդ որում $\varphi = 0$ ուղիղն անցնում է միավոր շրջանագծի գագաթների զույգով: Այսպիսով, $4n$-անկյուններով տրված նորմերն արդեն չեն համապատասխանում: Այնուհետև, այդ գագաթներում գոյություն ունի շոշափողների զույգ այնպիսին, որ դրանք շոշափում են միավոր շրջանագիծը միայն մեկ կետում, և դրանց հատման կետը պատկանում է $\varphi = -2$ ուղղին: Դիտարկենք գագաթները միացնող տրամագծով և հարևան հարթեցմամբ կազմված անկյունը: Նկարագրված պայմաններում այն պետք է փոքր լինի տրամագծի և շոշափողի միջև եղած անկյունից: Անցնենք դրանց տանգենսների անհավասարությանը.
\begin{equation*}
    \operatorname{tg}\frac{\pi(n-2)}{2n} < 2\cos\frac{\pi}{n}:
\end{equation*}
Այս անհավասարությունը տեղի չի ունենում, երբ $n \ge 6$: Այսինքն՝ ստանում ենք հակասություն բոլոր $\lambda$-հարթությունների համար: 

\textbf{Սահմանում 4.5.} Հարթության վրա դիտարկենք $ABC$ եռանկյունը, որի բոլոր անկյունները փոքր են $120^\circ$-ից: Այդ դեպքում $T$-ն \textit{Տորիչելիի կետն} է $A, B, C$ գագաթների համար, եթե $ATB, BTC, ATC$ անկյունները հավասար են $120^\circ$-ի:

\textbf{Դիտողություն 4.6.} Օգտագործենք $\hat{x}$ նշանակումը $\frac{1}{\|x\|}x$-ի համար, երբ $x \ne 0$:

\textbf{Պնդում 4.7.} \textit{Դիցուք $A = \{x_1, x_2, x_3\}$-ը $\lambda$-նորմավորված հարթության վրա մեկ ուղղի վրա չգտնվող կետերի բազմություն է, և $\lambda \equiv 0 \pmod 3$: Դիցուք $x_1, x_2, x_3$ գագաթներով եռանկյան մեջ բոլոր անկյունները փոքր են $120^\circ$-ից, և $p$-ն Տորիչելիի կետն է այդ եռանկյան համար: Այդ դեպքում, եթե $\widehat{x_i - p}$ կետերը չեն համընկնում $2\lambda$-անկյան գագաթների հետ, ապա $\operatorname{ft}(A)$-ն չվերասերված ուռուցիկ բազմանկյուն է, հակառակ դեպքում՝ $\operatorname{ft}(A) = \{p\}$:}

\textbf{Ապացույց:} Կառուցենք $\varphi_1, \varphi_2, \varphi_3$ նորմավորող ֆունկցիոնալները $x_1 - p, x_2 - p, x_3 - p$ կետերի համար: Եթե $\widehat{x_i - p}$-ը պատկանում են միավոր շրջանագծի հարթեցումներին, ապա ընդունենք $\varphi_i = 1$ այդ հարթեցման վրա: Եթե $\widehat{x_i - p}$-ը բազմանկյան գագաթներ են, ապա ընտրենք հարթեցումները շրջանցման միևնույն ուղղությամբ: Այդ դեպքում $\varphi = \varphi_1 + \varphi_2 + \varphi_3 = 0$ և ըստ Թեորեմ 2.4-ի՝ $p \in \operatorname{ft}(A)$: Ըստ Թեորեմ 2.6-ի՝ $\operatorname{ft}(A) = \bigcap_{i=1}^3 C(x_i, \varphi_i)$: Առաջին դեպքում յուրաքանչյուր $C(x_i, \varphi_i)$ կոն $p$-ի հետ միասին պարունակում է $p$-ի որոշակի շրջակայք, և $\operatorname{ft}(A)$-ն չվերասերված ուռուցիկ բազմանկյուն է, մյուս դեպքում $p$ կետն ընկած է $C(x_i, \varphi_i)$ կոներից յուրաքանչյուրի սահմանին, և համաչափության ուժով՝ $\bigcap_{i=1}^3 C(x_i, \varphi_i) = \{p\}$: 

\textbf{Թեորեմ 4.8.} \textit{$\lambda$-նորմավորված հարթության վրա Ֆերմա--Տորիչելիի խնդրի լուծումը միակ է ցանկացած երեք կետերի համար այն և միայն այն դեպքում, երբ $\lambda \not\equiv 0 \pmod 3$:}

\subsection{Կանոնավոր բազմանիստերը և եռաչափ նորմավորված տարածությունները}

[8] աշխատանքում ցույց էր տրվել չափանիշի կիրառությունը լյամբդա-հարթությունների՝ կանոնավոր բազմանկյուններով տրված նորմավորված հարթությունների վրա: Դիտարկենք խնդիրը որոշ եռաչափ տարածություններում:

\textbf{Պնդում 4.9.} \textit{Գոյություն ունի ընդամենը հինգ կանոնավոր բազմանիստ՝ տետրաէդր, խորանարդ, օկտաէդր, իկոսաէդր և դոդեկաէդր: Դրանցից չորսը բացի տետրաէդրից, որոշում են նորմ եռաչափ տարածությունում:}

Ստուգենք, արդյոք միա՞կն է Ֆերմա--Տորիչելիի խնդրի լուծումը ցանկացած երեք կետերի համար տվյալ տարածություններում: Դրա համար անհրաժեշտ է հետազոտել դրանցում գոյություն ունեցող նիստերի համաձայնեցված բազմությունները:

\textbf{Լեմմա 4.10.} \textit{Եթե Մինկովսկու տարածության միավոր սֆերայի երեք նիստերը կազմում են համաձայնեցված բազմություն, ապա դրանք զույգ առ զույգ չեն հատվում:}

\textbf{Ապացույց:} Համաձայնեցված բազմության յուրաքանչյուր նիստին համապատասխանում է ֆունկցիոնալ, որի մակարդակի մակերևույթը պարունակում է տվյալ նիստը: Բոլոր նիստերի ֆունկցիոնալների գումարը հավասար է զրոյի: Եթե ինչ-որ երկու նիստեր ունեն ընդհանուր կետ, ապա այդ կետում դրանց ֆունկցիոնալների գումարը հավասար է 2-ի: Այդ դեպքում երրորդ ֆունկցիոնալն այդ կետում ընդունում է $-2$ արժեքը, և այն չի հանդիսանում նորմավորող: Հակասություն: 

\textbf{Թեորեմ 4.11.} \textit{Կանոնավոր խորանարդով տրված նորմավորված տարածությունում գոյություն ունեն երեք կետեր, որոնց համար Ֆերմա--Տորիչելիի խնդրի լուծումը ոչ միակ է: Ավելին, եթե լուծումը ոչ միակ է, ապա այն հանդիսանում է հատված:}

\textbf{Ապացույց:} Դիցուք $X$-ը դիտարկվող տարածությունն է: Դիտարկենք $X$-ում համաձայնեցված բազմությունների բոլոր հնարավոր տարբերակները:

Դիցուք $X$-ում գոյություն ունի համաձայնեցված բազմություն, և դրանցից առնվազն երկուսը երկչափ նիստեր են: Ըստ լեմմա 4.10-ի՝ դրանք չունեն ընդհանուր կետեր: Ուրեմն, դրանք հակադիր են, և դրանց ֆունկցիոնալների գումարը հավասար է զրոյի: Բայց այդ դեպքում երրորդ ֆունկցիոնալը զրոյական է: Հակասություն:

Դիցուք համաձայնեցված բազմությունը պարունակում է ճիշտ մեկ երկչափ նիստ: Այն չի կարող ունենալ ընդհանուր կետեր մյուս երկու նիստերից որևէ մեկի հետ: Եթե մնացած երկու նիստերը խորանարդի կողեր են, ապա դրանք պատկանում են հակադիր երկչափ նիստին և հարևան չեն: Երկչափ նիստին համապատասխանող ֆունկցիոնալն ընդունում է $-1$ արժեքը երկու կողերի վրա էլ: Այդ դեպքում մյուս երկու ֆունկցիոնալներից յուրաքանչյուրը հավասար է զրոյի մյուս կողի վրա: Դա նշանակում է, որ դրանց մակարդակի հարթությունները խորանարդի հարևան նիստերի հետ կազմում են 45 աստիճանի անկյուն: Հարթությունները հատվում են մի ուղղով, որի վրա երկու ֆունկցիոնալների գումարը հավասար է 2-ի: Սակայն առաջին ֆունկցիոնալն այդ ուղղի վրա ընդունում է $-\frac{3}{2}$ արժեքը: Ֆունկցիոնալների գումարը հավասար չէ զրոյի: Հակասություն: Երկրորդ և երրորդ նիստերը չեն կարող լինել նաև կող և կետ կամ երկու կետ, քանի որ այդ դեպքում դրանց գծային թաղանթը չի պարունակում ուղիղ, ինչպես որ պետք է լինի չափանիշի պայմանում:

Դիցուք համաձայնեցված բազմությունը բաղկացած է երեք կողերից: Եթե ինչ-որ երկու կողեր պատկանում են միևնույն նիստին, ապա երրորդը պարտադիր պատկանում է հակադիրին: Դիցուք այն համաչափ է առաջին երկուսից մեկին: Այդ դեպքում խնդիրը հանգում է հարթության վրա նորմը որոշող քառակուսու երեք գագաթների համար ֆունկցիոնալների կառուցմանը: Այդպիսի ֆունկցիոնալներ գոյություն ունեն: Դիցուք համաչափ գագաթներին համապատասխանող ֆունկցիոնալների մակարդակի գծերը հատվում են քառակուսու անկյունագծի շարունակության վրա 2 նորմով կետում: Իսկ երրորդ ֆունկցիոնալի մակարդակի գիծն անցնում է մնացած գագաթով՝ ուղղահայաց այդ անկյունագծին: Այդ դեպքում ֆունկցիոնալների գումարը հավասար է զրոյի: Նմանապես կառուցվում են ֆունկցիոնալները խորանարդի կողերի համար: Կողերի գծային թաղանթների հատումն ուղիղ է: Այդ կողերի վրա վերցնելով երեք կամայական ներքին կետեր՝ լուծման մեջ կստանանք այդ ուղղի վրա ընկած ոչ դատարկ հատված:

Եթե երրորդ կողը դասավորված է ոչ զուգահեռ, ապա բոլոր կողերի գծային թաղանթների հատումը կետ է: Եթե կողերի ոչ մի զույգ չի պատկանում միևնույն նիստին, ապա դարձյալ կետ ենք ստանում դրանց գծային թաղանթների հատման մեջ: Չափանիշի պայմանը կատարված չէ:

Եթե համաձայնեցված բազմությունը բաղկացած է երկու կողից և կետից, ապա այդ կետը պետք է ընկած լինի կողերից յուրաքանչյուրի հետ մեկ երկչափ ենթատարածությունում: Այդ դեպքում կստանանք, որ կետը հանդիսանում է երկու կողերից մեկի ծայրակետը: Եթե համաձայնեցված բազմությունը բաղկացած է մեկ կողից և երկու գագաթներից, ապա գագաթներով անցնող ուղիղը խորանարդի անկյունագիծն է, և այն պետք է ընկած լինի կողի հետ միևնույն հարթության մեջ, բայց այդպիսի կող գոյություն չունի:

Այսպիսով, $X$-ում գոյություն ունեն միայն երեք զուգահեռ կողերից բաղկացած համաձայնեցված բազմություններ, իսկ Ֆերմա--Տորիչելիի խնդրի լուծումը կա՛մ հատված է, կա՛մ կետ: Պնդումն ապացուցված է: 

\textbf{Թեորեմ 4.12.} \textit{Կանոնավոր օկտաէդրով տրված նորմավորված տարածությունում Ֆերմա--Տորիչելիի խնդրի լուծումը միակ է ցանկացած երեք կետերի համար:}

\textbf{Ապացույց:} Նորմավորված տարածությունում, որտեղ միավոր սֆերան կանոնավոր օկտաէդր է, $a = (x, y, z)$ կոորդինատներով կետի նորմը տրվում է $\|a\| = |x| + |y| + |z|$ բանաձևով: Այդ դեպքում $a_i = (x_i, y_i, z_i)$ կետերի և $a = (x, y, z)$ կետի համար ունենք.
\begin{equation*}
    \sum_{i=1}^n |a a_i| = \sum_{i=1}^n \|a - a_i\| = \sum_{i=1}^n (|x - x_i| + |y - y_i| + |z - z_i|) = \sum_{i=1}^n |x - x_i| + \sum_{i=1}^n |y - y_i| + \sum_{i=1}^n |z - z_i|:
\end{equation*}
Այսինքն՝ խնդիրը հանգում է ըստ կոորդինատների լուծման, իսկ ըստ պնդում 2.8-ի՝ մեկ ուղղի վրա գտնվող երեք կետերի համար լուծումը միակ է: Պնդումն ապացուցված է: 

\textbf{Թեորեմ 4.13.} \textit{Կանոնավոր դոդեկաէդրով տրված նորմավորված տարածությունում գոյություն ունեն երեք կետեր, որոնց համար Ֆերմա--Տորիչելիի խնդրի լուծումը ոչ միակ է: Ավելին, եթե լուծումը ոչ միակ է, ապա այն հանդիսանում է հատված:}

\textbf{Ապացույց:} Դիտարկենք $X$-ում համաձայնեցված բազմությունների բոլոր հնարավոր տարբերակները: Հարմարության համար ներմուծենք նշանակումներ և ձևակերպումներ դոդեկաէդրի տարրերի համար: Երկչափ նիստերի նշանակման համար օգտագործվում են $f_i$-երը: $f_i, f_j$ նիստերին պատկանող կողը կնշանակենք $e_{i,j}$, իսկ $f_i, f_j, f_k$ նիստերին պատկանող գագաթը՝ $v_{i,j,k}$: $f_1$ և $f_{12}$ նիստերն անվանենք համապատասխանաբար վերին և ստորին: Կողերը և գագաթները պատկանում են դոդեկաէդրի հասարակածին, եթե չեն հատվում վերին և ստորին նիստերի հետ:

Ենթադրենք, որ համաձայնեցված բազմությունը պարունակում է առնվազն երկու երկչափ նիստ: Ըստ լեմմա 4.10-ի՝ դրանք չեն կարող լինել հարևան և հակադիր, ուստի կհամարենք, որ դրանք $f_1$ և $f_7$ նիստերն են: Եթե համաձայնեցված բազմության երրորդ տարրը նույնպես երկչափ նիստ է, ապա դա կարող է լինել միայն $f_9$-ը կամ $f_{10}$-ը: Քանի որ դրանք դասավորված են համաչափ առաջին երկու նիստերի նկատմամբ, ապա միաժամանակ լրացնում կամ չեն լրացնում դրանք մինչև համաձայնեցված բազմություն: Բայց եթե լրացնում են, ապա դրանց ֆունկցիոնալները պետք է համընկնեն, ինչը ճիշտ չէ:

Դիցուք երրորդ տարրը կող է: Բոլոր կողերը, բացի $e_{9,10}$-ից, ունեն դիտարկվող երկու նիստերի նկատմամբ համաչափ կողեր, ուստի չեն կարող մտնել համաձայնեցված բազմության մեջ: $e_{9,10}$ կողի ֆունկցիոնալը հավասար է $-1$-ի $v_{2,6,7}$ գագաթում: Նույն այդ գագաթում $f_7$ նիստի ֆունկցիոնալը հավասար է 1-ի, բայց $f_1$ նիստի ֆունկցիոնալը հավասար չէ զրոյի: Ֆունկցիոնալների գումարը հավասար չէ զրոյի:

Դիցուք երրորդ տարրը գագաթ է: Նախորդ դատողությունների նմանությամբ՝ բավական է ստուգել միայն $v_{4,9,10}$ գագաթը: Այս դեպքում $v_{2,6,7}$ գագաթում ֆունկցիոնալների գումարը դարձյալ հավասար չէ զրոյի:

Դիտարկենք այն տարբերակը, երբ համաձայնեցված բազմությունը պարունակում է միակ $f_1$ նիստը: Դիցուք մյուս երկու տարրերը կողեր են: Եթե կողերից մեկը պատկանում է ստորին նիստին, ապա նրանով անցնում է երկրորդ կողի ֆունկցիոնալի զրոյական մակարդակը: Սակայն այս դեպքում դա վերին նիստի հետ սահմանակից կողերից մեկն է, ինչն անհնար է ըստ լեմմա 4.10-ի: Դիցուք կողերից մեկը միացնում է հասարակածը և ստորին նիստը: Բայց այդ դեպքում մնացած կողերից ցանկացածն ունի դիտարկվող նիստի և կողի նկատմամբ համաչափ կող և չի կարող դրանք լրացնել մինչև համաձայնեցված բազմություն: Այդ դեպքում երկու կողերն էլ պետք է պատկանեն հասարակածին, ընդ որում դրանք համաչափ են

\begin{figure}[htbp]
\centering
\includegraphics[width=0.45\textwidth]{images/5.png}
\caption{Կանոնավոր դոդեկաէդրը որոշում է նորմ եռաչափ տարածությունում}
\label{fig:5}
\end{figure}

կենտրոնի նկատմամբ: Այդ կողերով անցնող հարթությունը $f_1$ նիստի ֆունկցիոնալի զրոյական մակարդակն է, բայց զուգահեռ չէ նրան, ինչն անհնար է:

Դիցուք $f_1$ նիստը լրացնում են միևնույն հարթության մեջ ընկած կողը և գագաթը: Հասարակածի կողերով տրված ենթատարածություններում ընկած են միայն այն գագաթները, որոնք հանդիսանում են բուն կողի և նրան հակադիր կողի ծայրակետերը: Դրանցից յուրաքանչյուրն ունի իրեն համաչափ գագաթ: Նմանապես ստորին նիստի կողը նույնպես չի կարող մտնել համաձայնեցված բազմության մեջ: Մնացին ստորին նիստի և հասարակածի միջև եղած կողերը: Բավական է ստուգել $e_{8,9}$ կողը և $v_{5,6,11}$ գագաթը: Այս դեպքում $v_{5,6,11}$ և $v_{3,8,9}$ գագաթները միացնող ուղղի վրա $f_1$ նիստի ֆունկցիոնալը հավասար է զրոյի: Բայց այդ ուղիղը զուգահեռ չէ նիստի հարթությանը: Հակասություն:

Եթե $f_1$ նիստը լրացնում են երկու համաչափ գագաթներ, ապա դա անպայման հասարակածի գագաթների զույգ է: Բայց դրանցով անցնող ուղիղը զուգահեռ չէ նիստի հարթությանը: Հակասություն, քանի որ նիստի ֆունկցիոնալն այդ ուղղի վրա պետք է հավասար լինի զրոյի:

Մնացին համաձայնեցված բազմությունների այն տարբերակները, որոնք չեն պարունակում և ոչ մի նիստ: Ենթադրենք, որ համաձայնեցված բազմությունը բաղկացած է երեք կողերից, որոնց ենթատարածությունների հատումը պարունակում է առնվազն ուղիղ: Դիցուք դրանցից երկուսը հակադիր են, օրինակ՝ դրանք $e_{1,4}$-ն ու $e_{7,12}$-ն են: Դրանց ընդհանուր երկչափ ենթատարածությունը որոշում է երրորդ կողի ֆունկցիոնալի զրոյական մակարդակը: Տվյալ հարթությանը զուգահեռ են միայն $e_{2,6}$ և $e_{9,10}$ կողերը: Դրանք համաչափ են, ուստի բավական է դիտարկել առաջինը:

Դիցուք $\varphi_1$ ֆունկցիոնալի մակարդակն անցնում է $e_{2,6}$ կողով և հավասար երկնիստ անկյուններ է կազմում հարևան նիստերի հետ: $\varphi_2$ և $\varphi_3$ ֆունկցիոնալների մակարդակները պարունակում են համապատասխանաբար $e_{1,4}$ և $e_{7,12}$ կողերը: Այդ երկու կողերի հարթությունում ֆունկցիոնալների գումարը հավասար է զրոյի: Այժմ դիտարկենք $e_{2,6}$ կողը և $e_{1,4}$ ու $e_{7,12}$ կողերի միջնակետերը պարունակող հարթությունը: Դոդեկաէդրի հատույթն այդ հարթությամբ պատկերված է նկար 6-ում: Քանի որ ֆունկցիոնալների գումարը հավասար է զրոյի, ապա $\varphi_2 = 1$ և $\varphi_3 = 1$ մակարդակի գծերի հատման կետը պատկանում է $\varphi_1 = -2$ ուղղին: Դիցուք $\alpha$-ն երկնիստ անկյունն է $\varphi_1 = 0$ մակարդակի և նիստի միջև, իսկ $\beta$-ն՝ անկյունը $\varphi_1 = 0$ և $\varphi_2 = 1$ մակարդակների միջև: Եթե $\beta > \alpha$, ապա այդպիսի ֆունկցիոնալներ գոյություն ունեն, և համապատասխան կողերի բազմություն նույնպես:

\begin{figure}[htbp]
\centering
\includegraphics[width=0.45\textwidth]{images/6.png}
\caption{}
\label{fig:6}
\end{figure}

Ապացուցենք, որ դա իսկապես այդպես է: $\alpha$ անկյունը կանոնավոր դոդեկաէդրի երկնիստ անկյան կեսն է, այսինքն՝
\begin{equation*}
    \alpha = \frac{1}{2}\arccos\left(-\frac{1}{\sqrt{5}}\right):
\end{equation*}
$e_{2,6}$ և $e_{9,10}$ կողերի երկարությունները հավասար են $a = 1$: Հատույթում ստացված վեցանկյան մնացած կողմերի երկարությունները հավասար են կանոնավոր հնգանկյան բարձրությանը, այսինքն՝ $b = \frac{\sqrt{5+2\sqrt{5}}}{2}$: Հաշվենք $b$ ներքնաձիգով և $\alpha$ սուր անկյունով ուղղանկյուն եռանկյան էջերը.
\begin{equation*}
    c = b\cos\alpha = \frac{1}{2}\sqrt{\frac{3+\sqrt{5}}{2}}, \quad d = b\sin\alpha = \frac{1}{2}\sqrt{\frac{7+3\sqrt{5}}{2}}:
\end{equation*}
Այժմ կարող ենք հաշվել $\beta$ անկյան տանգենսը.
\begin{equation*}
    \operatorname{tg}\beta = \frac{2d}{c + \frac{a}{2}} = \frac{\sqrt{\frac{7+3\sqrt{5}}{2}}}{\frac{1}{2}\sqrt{\frac{3+\sqrt{5}}{2}} + \frac{1}{2}} = 2:
\end{equation*}
Քանի որ $\operatorname{tg} 2 > \frac{1}{2}\arccos\left(-\frac{1}{\sqrt{5}}\right)$, ապա $\beta > \alpha$, և դոդեկաէդրի կողերը կազմում են համաձայնեցված բազմություն: Քանի որ այդ երեք կողերով տրված երկչափ ենթատարածությունները հատվում են ուղղով, ապա տվյալ բազմությունը տալիս է Ֆերմա--Տորիչելիի խնդրի ոչ միակ լուծումներ, ընդ որում՝ որոշակի ուղղին պատկանող, այսինքն՝ հատվածներ:

Եթե դոդեկաէդրի երեք կողերի մեջ չկա հակադիր կողերի ոչ մի զույգ, ապա կողերը պարունակող երկչափ ենթատարածությունները հատվում են ոչ ավելի, քան կետով:

Դոդեկաէդրի համաձայնեցված բազմությունը չի կարող բաղկացած լինել երկու կողից և գագաթից, քանի որ այդ դեպքում գագաթը պետք է ընկած լինի կողերից յուրաքանչյուրի հետ միևնույն հարթության մեջ: Ստանում ենք, որ կողերը հակադիր են, իսկ գագաթը կողերից մեկի ծայրակետն է, ինչն անհնար է ըստ լեմմա 4.10-ի:

Նույնպես գոյություն չունի միևնույն հարթության մեջ ընկած մեկ կողից և երկու գագաթներից բաղկացած համաձայնեցված բազմություն: Եթե դա այդպես լիներ, ապա ըստ լեմմա 4.10-ի՝ գագաթները կլինեին տվյալ կողին հանդիպակաց կողի ծայրակետերը: Օգտագործելով նկար 6-ը՝ ենթադրենք, որ դիտարկվում է $e_{2,6}$ կողը, իսկ գագաթները $e_{9,10}$ կողի ծայրակետերն են: Այդ դեպքում $\varphi_2$ և $\varphi_3$ ֆունկցիոնալների մակարդակի գծերն անցնում են տվյալ գագաթներով և հատվում $\varphi_1 = -2$ գծի վրա: Նկարից երևում է, որ դա անհնար է:

Այսպիսով, դոդեկաէդրն ունի միայն երեք կողերից բաղկացած համաձայնեցված բազմություններ: Ընդ որում, եթե դրանց օգնությամբ ստացվել է ոչ միակ լուծում, ապա դա հատված է: Պնդումն ապացուցված է: 

\textbf{Թեորեմ 4.14.} \textit{Եթե եռաչափ նորմավորված տարածության միավոր սֆերան պրիզմա է, ապա ցանկացած $n \ge 2$-ի համար գոյություն ունեն $n$ կետեր, որոնց համար Ֆերմա--Տորիչելիի խնդրի լուծումը ոչ միակ է:}

\textbf{Ապացույց:} Դիտարկենք նորմավորված հարթություն, որի միավոր շրջանագիծը համընկնում է պրիզմայի հիմքի հետ: Այդ նորմավորված հարթությունում վերցնենք ցանկացած համաձայնեցված բազմություն և այն որոշող ֆունկցիոնալները: Տանենք այդ ֆունկցիոնալների մակարդակի գծերը պրիզմայի երկու հիմքերի մոտ: Դրանցով անցնում է հենման հարթությունների միակ հավաքածուն, որը որոշում է $n$ ֆունկցիոնալ ելակետային տարածությունում: Դրանց գումարը հավասար է զրոյի, իսկ նիստերի գծային թաղանթը պարունակում է ուղիղ: 3.9 չափանիշի պայմանը կատարված է, և պնդումն ապացուցված է: 

\section{Կոնֆիգուրացիոն տարածությունները և լուծումների կայունությունը}

Դիցուք $X$-ը $d$ չափայնության նորմավորված տարածություն է: Դիտարկենք նրանում Ֆերմա--Տորիչելիի խնդիրը $n$ կետերի համար:

\textbf{Սահմանում 5.1.} $C_{X,n}$ \textit{կոնֆիգուրացիոն տարածություն} է կոչվում $\mathbb{R}^{nd}(x_{11}, \dots, x_{1d}, \dots, x_{n1}, \dots, x_{nd})$ էվկլիդյան տարածությունը, որի կետերին համապատասխանում են $X$ տարածության կետերի հավաքածուները:

\subsection{Ֆերմա--Տորիչելիի խնդրի լուծումների կայունությունը}

Նորմավորված տարածության կետերի հավաքածուները բաժանվում են երկու դասի՝ կախված Ֆերմա--Տորիչելիի խնդրի դրանց լուծման միակության բնույթից: Որպեսզի ուսումնասիրենք, թե ինչպես է կառուցված դրանց միջև սահմանը կոնֆիգուրացիոն տարածությունում, դիտարկենք լուծումների կայունության հատկությունը:

\textbf{Սահմանում 5.2.} Դիցուք Ֆերմա--Տորիչելիի խնդրի լուծումը $x = (x_1, \dots, x_n)$ կետերի հավաքածուի համար միակ է (ոչ միակ է): Կասենք, որ լուծումը \textit{կայուն է}, եթե գոյություն ունի $x$ հավաքածուի այնպիսի շրջակայք $C_{X,n}$ կոնֆիգուրացիոն տարածությունում, որ այդ շրջակայքի ցանկացած հավաքածուի համար Ֆերմա--Տորիչելիի խնդրի լուծումը նույնպես միակ է (ոչ միակ է):

$X$-ն անվանենք \textit{կայուն $n$-միակությամբ ($n$-ոչ միակությամբ) տարածություն}, եթե $n$ կետերի համար ցանկացած միակ (ոչ միակ) լուծում հանդիսանում է կայուն:

Նախ ցույց տանք, որ լուծման միակության կայունության հարցը տրիվիալ է:

\textbf{Պնդում 5.3.} \textit{Դիցուք $n$-ը կենտ է: $X$ տարածությունը հանդիսանում է կայուն $n$-միակությամբ տարածություն այն և միայն այն դեպքում, երբ Ֆերմա--Տորիչելիի խնդրի լուծումը միակ է ցանկացած $n$ կետերի համար:}

\textbf{Ապացույց:} Ենթադրենք, որ $X$-ում կա կետերի այնպիսի հավաքածու, որ լուծումը ոչ միակ է: Կարելի է համարել, որ այն պարունակում է զրո: Դիտարկենք այդ լուծմանը համապատասխանող նիստերի համաձայնեցված բազմությունը: Դիցուք $x = (x_1, \dots, x_n)$-ը այդ նիստերի ներքին կետերի հավաքածուն է: Ցանկացած $\lambda > 0$-ի համար $\lambda x = (\lambda x_1, \dots, \lambda x_n)$ հավաքածուի համար խնդրի լուծումը նույնպես ոչ միակ է: Այդ դեպքում զրոյի ցանկացած շրջակայքում ընկած է ոչ միակ լուծումով $\lambda x$ հավաքածու բավականաչափ փոքր $\lambda > 0$-ի դեպքում: Ընդ որում բուն զրոյական հավաքածուի համար լուծումը միակ է: 

\textbf{Պնդում 5.4.} \textit{Զույգ $n$-ի դեպքում կայուն $n$-միակությամբ տարածություններ գոյություն չունեն:}

\textbf{Ապացույց:} Մեկ ուղղի վրա ընկած տարբեր կետերի զույգ քանակի համար լուծումը միշտ ոչ միակ է: Նախորդ ապացուցման նմանությամբ ստանում ենք, որ զրոյական հավաքածուն անկայուն է: 

Ոչ միակ լուծումների կառուցվածքի ուսումնասիրությունն սկսենք վեցանկյուն նորմավորված հարթությունների դիտարկումից:

\textbf{Լեմմա 5.5.} \textit{Յուրաքանչյուր վեցանկյուն նորմավորված հարթության համար նրա նիստերի համաձայնեցված բազմությունների հավաքածուն պատկանում է երեք տիպերից մեկին.}
\begin{enumerate}
    \item \textit{Գոյություն ունեն միայն երեք նիստ-հատվածներից և երեք նիստ-կետերից բաղկացած համաձայնեցված բազմություններ $\iff$ Ոչ միակ լուծումներ են հանդիսանում միայն բազմանկյունները:}
    \item \textit{Գոյություն ունեն բոլոր տեսակի համաձայնեցված բազմությունները, բացի երեք նիստ-հատվածներից $\iff$ Ոչ միակ լուծումներ են հանդիսանում միայն հատվածները:}
    \item \textit{Գոյություն ունեն միայն երեք նիստ-կետերից բաղկացած համաձայնեցված բազմություններ $\iff$ Ոչ միակ լուծումներ գոյություն չունեն:}
\end{enumerate}

\textbf{Ապացույց:} Ենթադրենք, որ վեցանկյուն նորմավորված հարթության երեք նիստ-հատվածները կազմում են համաձայնեցված բազմություն: Վեցանկյան երկու ոչ հարևան նիստերի վրա հենվող երկու ֆունկցիոնալների համար դրանք մինչև զրո լրացնող ֆունկցիոնալը հենվում է երրորդ նիստի վրա, ուստի երկու հատվածից և կետից բաղկացած համաձայնեցված բազմություն գոյություն չունի: Եթե համաձայնեցված բազմության մեջ մտնում են երկու կետեր, ապա դրանք հակադիր են: Դրանց վրա հենվող ֆունկցիոնալները չեն կարող նիստի ֆունկցիոնալը լրացնել մինչև զրո, քանի որ հենման ուղիղների հատման կետն ընկած է դրանց հարակից նիստերի հենման ուղիղների հատման կետից ավելի հեռու, որոնք իսկապես առաջին ֆունկցիոնալը լրացնում են մինչև զրո:

Դիցուք գոյություն ունի երկու նիստ-կետերից և նիստ-հատվածից բաղկացած համաձայնեցված բազմություն: Նիստ-կետերը հակադիր են, դրանցով անցնում է նիստ-հատվածի ֆունկցիոնալի զրոյական մակարդակը: Կարելի է նիստ-կետերի ֆունկցիոնալները պտտել բուն կետերի շուրջն այնպես, որ դրանց հատման կետը մնա առաջին ֆունկցիոնալի $\varphi = -2$ մակարդակի գծի վրա, ընդ որում ֆունկցիոնալների գումարը կմնա զրոյական: Պտտենք դրանք այնպես, որ ֆունկցիոնալներից մեկը դառնա հենման ֆունկցիոնալ հենման կետին հարակից նիստի համար: Կստանանք երկու նիստ-հատվածներից և մեկ նիստ-կետից բաղկացած համաձայնեցված բազմություն: Նմանապես կարելի է համաձայնեցված երկու հատվածից և կետից ստանալ համաձայնեցված երկու կետ և մեկ հատված: Պնդումն ապացուցված է: 

Ինչպես ցույց կտրվի հետագայում, վեցանկյուն հարթության երկու տիպերն էլ օժտված են կայուն 3-ոչ միակությամբ: Նկար 7-ում այս փաստը պատկերված է բիֆուրկացիոն դիագրամների տեսքով:

\begin{figure}[htbp]
\centering
\includegraphics[width=0.75\textwidth]{images/7.png}
\caption{Ոչ միակ լուծումները վեցանկյուն նորմերում}
\label{fig:7}
\end{figure}

\textbf{Լեմմա 5.6.} \textit{Դիցուք $x = (x_1, \dots, x_n)$ հավաքածուի համար լուծումը միակ է (ոչ միակ է) և գոյություն ունեն շրջակայքեր բոլոր $x_i$ կետերի համար, բացի մեկ $x_n$-ից, այնպիսին, որ ցանկացած $x' = (x'_1, \dots, x'_{n-1}, x_n)$ հավաքածուի համար, որտեղ $x'_i$-ն $x_i$-ի շրջակայքից է, լուծումը նույնպես միակ է (ոչ միակ է): Այդ դեպքում լուծումը կայուն է:}

\textbf{Ապացույց:} Ըստ լուծման կայունության սահմանման՝ միակությունը կամ ոչ միակությունը պետք է պահպանվի հավաքածուի որոշակի շրջակայքի բոլոր կետերի համար: Դա համարժեք է լուծման պահպանմանն այն հավաքածուների համար, որտեղ յուրաքանչյուր բաղադրիչ վերցված է ելակետային հավաքածուի բաղադրիչի շրջակայքից: Սակայն, քանի որ կետերի հավաքածուի զուգահեռ տեղափոխությունը պահպանում է Ֆերմա--Տորիչելիի խնդրի լուծումը, ապա կարելի է սահմանափակվել այն հավաքածուներով, որտեղ բաղադրիչներից մեկը մնում է տեղում: 

\textbf{Պնդում 5.7.} \textit{Բոլոր նորմավորված հարթությունները, որոնց վրա նորմը տրվում է բազմանկյուններով, օժտված են կայուն 3-ոչ միակությամբ:}

\textbf{Ապացույց:} Դիցուք $x = (x_1, x_2, x_3)$ կետերի համար Ֆերմա--Տորիչելիի խնդրի լուծումը ոչ միակ է, և զրոն մտնում է լուծման մեջ: $\varphi_1, \varphi_2, \varphi_3$ ֆունկցիոնալները նորմավորող են $x_1, x_2, x_3$-ի համար: Անհրաժեշտ է դիտարկել $\varphi_i$ հավաքածուին համապատասխանող նիստերի համաձայնեցված բազմության կազմի երեք տարբերակ:

\textbf{Դեպք 1:} Դիցուք յուրաքանչյուր $\varphi_i$ ֆունկցիոնալ հենվում է նիստ-հատվածի վրա:

Եթե գոնե երկու $\hat{x}_1, \hat{x}_2$ կետեր պատկանում են հարթեցումների ներքին տիրույթներին, ապա դրանց համար գոյություն ունեն այնպիսի շրջակայքեր, որ ցանկացած $x'_1, x'_2, x_3$ հավաքածուի համար լուծումը ոչ միակ է: Այդ դեպքում, ըստ լեմմա 5.6-ի, լուծումը կայուն է:

Եթե երկու $\hat{x}_i$ և $\hat{x}_j$ կետերը միավոր շրջանագծի ոչ հակադիր գագաթներ են, ապա զուգահեռ տեղափոխությամբ հավաքածուն շեղենք միավոր շրջանագծի այն անկյունագծի երկայնքով, որի վրա ընկած է $\hat{x}_i$ կետը: Տեղափոխության ուղղությունն անհրաժեշտ է տալ այնպես, որ $x_j$ կետը պահպանի իր նորմավորող ֆունկցիոնալը: Ստացանք նախորդ դեպքը, որի համար լուծումը կայուն է:

Եթե $\hat{x}_i$ և $\hat{x}_j$ կետերը հակադիր են, ապա այն նիստերը, որոնց վրա հենվում են դրանց նորմավորող ֆունկցիոնալները, ընկած են $x_i x_j$ հատվածը պարունակող ուղղի միևնույն կողմում: Ուրեմն, կարելի է զուգահեռ տեղափոխությամբ $x_1, x_2, x_3$ հավաքածուն շեղել $x_i x_j$ հատվածին ուղղահայաց ուղղությամբ այնպես, որ $\hat{x}_i, \hat{x}_j$ կետերն ընկնեն հարթեցումների ներսը:

Դիցուք բոլոր $\hat{x}_1, \hat{x}_2, \hat{x}_3$ կետերը չեն ընկած հարթեցումների ներսում: Ստուգենք շեղումները միավոր շրջանագծի այն տրամագծերի երկայնքով, որոնք պարունակում են $x_i$ կետերը: Եթե որևէ տրամագծի համար գոյություն ունի շեղում, որից հետո մյուս երկու կետերն ընկնում են իրենց ֆունկցիոնալներին համապատասխանող հարթեցումների ներսը, ապա ստանում ենք կայունություն: Եթե այդպիսի շեղում գոյություն չունի, ապա $x_i$-ից ելնող կոները հատվում են միայն զրոյով, և լուծումը միակ է:

\textbf{Դեպք 2:} Դիցուք $\varphi_1$ և $\varphi_2$ ֆունկցիոնալները հենվում են հարթեցումների վրա, իսկ $\varphi_3$-ը՝ գագաթի: Եթե $\hat{x}_1, \hat{x}_2$ կետերն ընկած են հարթեցումների ներսում, ապա դրանց համար գոյություն ունեն այնպիսի շրջակայքեր, որ ցանկացած $x'_1, x'_2, x_3$ հավաքածուի համար լուծումը ոչ միակ է: Այդ դեպքում, ըստ լեմմա 5.6-ի, լուծումը կայուն է:

Եթե $\hat{x}_1, \hat{x}_2$ կետերից մեկը գագաթ է, ապա $x$ հավաքածուն շեղենք միավոր շրջանագծի $\hat{x}_3$ կետը պարունակող տրամագծի երկայնքով այնպես, որ բոլոր $\hat{x}_i$ կետերը պահպանեն իրենց ֆունկցիոնալները:

Եթե երկու $\hat{x}_1, \hat{x}_2$ կետերն էլ հանդիսանում են գագաթներ, ապա դարձյալ դիտարկենք $\hat{x}_3$-ը պարունակող տրամագծի երկայնքով շեղման հնարավորությունը: Այնպիսի շեղման գոյության դեպքում, որ $\hat{x}_1, \hat{x}_2$ կետերն ընկնում են հարթեցումների ներսը, ստանում ենք կայունություն: Եթե անկախ շեղման ուղղությունից կետերից մեկը չի պահպանում իր նորմավորող ֆունկցիոնալը, ապա ելակետային $\hat{x}_1, \hat{x}_2$ կետերը դասավորված են այնպես, որ դրանցից ելնող կոները $\hat{x}_3$-ը պարունակող տրամագծի վրա հատվում են միայն զրոյով: Այդ դեպքում լուծումը միակ է, և, հետևաբար, այդպիսի դեպք չի առաջանում:

\textbf{Դեպք 3:} Դիցուք $\varphi_1$ ֆունկցիոնալը հենվում է հարթեցման վրա, իսկ $\varphi_2$-ը և $\varphi_3$-ը՝ գագաթների: Այդ դեպքում այդ գագաթները հակադիր են, և կարելի է $\varphi_2$-ը և $\varphi_3$-ը պտտել շոշափման կետերի շուրջն այնպես, որ դրանցից մեկը սկսի հենվել հարթեցման վրա: Այսպիսով, միևնույն $x$ հավաքածուին այժմ համապատասխանում է երկու հարթեցումներից և գագաթից բաղկացած համաձայնեցված բազմություն: Այս դեպքն արդեն դիտարկված է, ուրեմն՝ լուծումը կայուն է: 

Կայուն 3-ոչ միակությամբ նորմերի լրիվ նկարագրության համար լրացնենք միավոր սֆերայի նիստերի համաձայնեցված բազմության սահմանումը [8]-ից:

\textbf{Սահմանում 5.8.} Վերցնենք Մինկովսկու տարածության $S$ միավոր սֆերայի $F$ նիստերի վերջավոր բազմություն: Յուրաքանչյուր $f_i \in F$-ի համար ընտրենք $\pi_i$ հենման հիպերհարթություն այնպես, որ $\pi_i \cap S = f_i$: Դիցուք $\pi_i$ հիպերհարթությունը որոշում է որոշակի $\varphi_i$ գծային ֆունկցիոնալի $\varphi_i = 1$ մակարդակի մակերևույթը: Եթե գտնվի հենման հիպերհարթությունների այնպիսի հավաքածու, որ կառուցված ֆունկցիոնալների գումարը հավասար լինի զրոյականին, ապա նիստերի այդպիսի բազմությունը կոչվում է \textit{համաձայնեցված}:

Նիստերի $F_0 = \{f_1^0, \dots, f_n^0\}$ համաձայնեցված բազմությունը կոչվում է \textit{մաքսիմալ}, եթե գոյություն չունի այլ $F = \{f_1, \dots, f_n\}$ համաձայնեցված բազմություն, որն ամբողջությամբ ծածկում է այն, այսինքն՝ այնպիսին, որ $f_i^0 \subset f_i$ ցանկացած $i = 1, \dots, n$-ի համար:

\textbf{Թեորեմ 5.9.} \textit{Նորմավորված հարթությունում գոյություն ունի անկայուն ոչ միակ լուծումով երեք կետերի հավաքածու այն և միայն այն դեպքում, երբ նրա միավոր շրջանագծի մոտ գտնվի երկու կետից և հարթեցումից բաղկացած նիստերի մաքսիմալ համաձայնեցված բազմություն:}

\textbf{Ապացույց:} 5.7 պնդման ապացուցման դատողությունը 1-ին և 2-րդ դեպքերի համար ճիշտ է նաև կամայական նորմավորված հարթությունների համար: Դեպք 3-ը նույնպես տեղի ունի, եթե համաձայնեցված հարթեցումից և երկու գագաթներից բաղկացած բազմությունը մաքսիմալ չէ: Այսպիսով, եթե որոնելի համաձայնեցված բազմությունը բացակայում է, ապա բոլոր ոչ միակ լուծումները կայուն են:

Ենթադրենք, որ միավոր շրջանագիծն ունի հարթեցումից և երկու գագաթներից բաղկացած $F_M$ համաձայնեցված բազմություն, և այն մաքսիմալ է: Այդ դեպքում այդ գագաթների որոշակի շրջակայքում բացակայում են հարթեցումները, ինչպես նաև այլ համաձայնեցված բազմությունների որևէ նիստեր: Դիցուք այդ շրջակայքին համապատասխանում է $\varepsilon$ բացվածքի անկյունով և զրոյում գագաթով կոնը: Դիտարկենք եռանկյան գագաթների հավաքածու, որի հիմքը զուգահեռ է $F_M$-ի գագաթները միացնող ուղղին, իսկ երրորդ գագաթը հավասարահեռ է մյուս երկու գագաթներից և գտնվում է հիմքից բավականաչափ փոքր հեռավորության վրա այնպես, որ եռանկյան սուր անկյունները փոքր են $\varepsilon$-ից: Այս հավաքածուի համար լուծումը ոչ միակ է և համապատասխանում է դիտարկվող $F_M$ նիստերին:

Փոքր հեռավորությամբ շեղենք եռանկյան հիմքի գագաթներից մեկն այնպես, որ եռանկյան հիմքը փոխի ուղղությունը: Ենթադրենք, որ նոր հավաքածուի լուծումը դարձյալ ոչ միակ է, այդ դեպքում լուծումներից մեկն ընկած է եռանկյան մեջ և չի հանդիսանում նրա գագաթներից մեկը: Եռանկյունը շեղենք այնպես, որ այդ լուծումը համընկնի զրոյի հետ: Զրոյից դեպի եռանկյան հիմքի գագաթները տարված ճառագայթները պետք է հատեն միավոր շրջանագծի համաձայնեցված նիստերը: Քանի որ այդ հատումները գտնվում են $F_M$-ի գագաթների $\varepsilon$-շրջակայքում, բայց դրանք չեն հանդիսանում, ապա ստանում ենք հակասություն: 

\subsection{Լուծումների բիֆուրկացիոն դիագրամները}

Դիտարկենք Ֆերմա--Տորիչելիի խնդիրը երեք կետերի համար նորմավորված հարթությունում: Գոյություն ունեն հարթության ձևափոխություններ, որոնք ցանկացած $x = (x_1, x_2, x_3)$ հավաքածուի համար դրանց $\operatorname{ft}(x)$ լուծումների բազմությունը փոխադրում են $\operatorname{ft}(x')$՝ փոփոխված հավաքածուի լուծումների բազմության: Այդպիսի ձևափոխություններ են հանդիսանում շեղումները և զրոյում կենտրոնով հոմոթետիաները:

Համարենք, որ երեք կետերի $x = (x_1, x_2, x_3)$ հավաքածուն տրվում է $x_1 x_2$ վեկտորի բևեռային $\varphi$ անկյունով և այդ վեկտորի նկատմամբ $x_3$ կետի դիրքով: Ֆիքսված $\varphi$ անկյան համար դիտարկենք լուծումների բիֆուրկացիոն դիագրամը՝ հարթությունը տրոհենք երկու ենթաբազմության: Դրանցից մեկը պատասխանում է $x_3$ կետի այն դիրքերին, որոնց համար $\operatorname{ft}(x)$ բազմությունը բաղկացած է միակ լուծումից, իսկ մյուսը՝ ոչ միակ լուծումներին համապատասխանող դիրքերին:

Ուսումնասիրենք բիֆուրկացիոն դիագրամների հատկությունները: Քանի որ մեկ ուղղի վրա ընկած ցանկացած երեք կետերի համար լուծումը միակ է և համընկնում է այդ կետերից մեկի հետ, ապա ստանում ենք առաջին հատկությունը.

\textbf{Պնդում 5.10.} \textit{$x_1$ և $x_2$ կետերը պարունակող ուղիղը մտնում է միակ լուծումների ենթաբազմության մեջ:}

Ֆերմա--Տորիչելիի խնդրի լուծման տեսքը կախված է միայն այն կոների սահմանների ուղղությունից, որոնց հատումն է այն հանդիսանում, որտեղից ճշմարիտ է հետևյալ պնդումը:

\textbf{Պնդում 5.11.} \textit{Միավոր շրջանագծի հարթեցումների յուրաքանչյուր մաքսիմալ համաձայնեցված բազմությանը համապատասխանեցնենք $\{(\varphi_{11}, \varphi_{12}), \dots, (\varphi_{n1}, \varphi_{n2})\}$ անկյունների հավաքածուն, որտեղ $(\varphi_{i1}, \varphi_{i2})$-ը $f_i$ նիստի ծայրակետերի բևեռային անկյուններն են: Եթե նիստը կետ է, ապա $\varphi_{i1} = \varphi_{i2}$: Անկյունների այս հավաքածուները միարժեքորեն որոշում են բիֆուրկացիոն դիագրամների տեսքը տվյալ նորմավորված հարթության համար:}

\textbf{Պնդում 5.12.} \textit{Դիցուք նորմավորված հարթության համար հայտնի է երեք համաձայնեցված նիստերից բաղկացած բոլոր մաքսիմալ հավաքածուների բազմությունը: Այդ դեպքում ցանկացած բիֆուրկացիոն դիագրամի համար ոչ միակ լուծումների ենթաբազմությունը տրոհվում է չհատվող տիրույթների այնպես, որ յուրաքանչյուր հավաքածուին համապատասխանում է մեկ տիրույթ:}

\begin{figure}[htbp]
\centering
\includegraphics[width=0.75\textwidth]{images/8.png}
\caption{Բիֆուրկացիոն դիագրամներ բազմանկյուն նորմի և հարթեցումներով նորմի համար}
\label{fig:8}
\end{figure}

\textbf{Ապացույց:} Երկու տարբեր տիրույթների հատումը կնշանակեր, որ կետերի միևնույն հավաքածուի համար կարելի է ընտրել երկու տարբեր մաքսիմալ համաձայնեցված բազմություններ, որոնք տալիս են համընկնող լուծում:

Ենթադրենք, որ դա այդպես է: Քանի որ հավաքածուները մաքսիմալ են, ապա դրանք տարբերվում են առնվազն երկու տարրով: Եթե մի հավաքածուի գոնե մեկ տարր չի հատվում մյուսի տարրերից ոչ մեկի հետ, ապա համապատասխան լուծումները նույնպես չեն հատվում: Դիցուք երկու տարրերը տարբեր են և հատվում են մյուս հավաքածուի տարրերի հետ կետով, այդ դեպքում կոները հատվում են ճառագայթով: Ստանում ենք, որ լուծումը հատված է: Բայց այս դեպքում բազմությունները բաղկացած են երկու նիստ-հատվածից և մեկ գագաթից, ընդ որում տարբեր բազմությունների գագաթները հակադիր են: Այդ դեպքում դրանք կարող են լինել լուծում կետերի միևնույն հավաքածուի համար միայն այն դեպքում, եթե երկու $x_1, x_2$ կետերն ընկած են գագաթները միացնող ուղղի վրա: Նիստերի երկու բազմություններն էլ կազմում են ոչ միակության միևնույն տիրույթը՝ $x_1 x_2$ հատվածի կետերից ելնող կոների միավորումը:

Այսպիսով, եթե ոչ միակության երկու տարբեր տիրույթներ հատվում են, ապա դրանք համընկնում են, ինչն էլ պահանջվում էր ապացուցել: 

\section{Եզրակացություն}

Ֆերմա--Տորիչելիի դասական խնդրի ձևակերպման պարզությունը թույլ է տալիս դիտարկել այս պրոբլեմի ընդհանրացման բազմաթիվ տարբերակներ: Իմ աշխատանքում ուսումնասիրվել է խնդրի ընդլայնումը էվկլիդյան հարթությունից կամայական վերջավորաչափ նորմավորված տարածությունների վրա: Լուծումների ոչ միակության պայմանների մասին ստացված արդյունքներն օգնում են վերլուծել նորմերի ամբողջ ընտանիքներ: Մասնավորապես, պատասխաններ են ստացվել լյամբդա-հարթությունների, կանոնավոր բազմանիստերի և որոշ այլ նորմերի համար:

Կոնֆիգուրացիոն տարածությունների հետազոտումը շարունակում է նորմավորված տարածությունների դասակարգման ընդլայնումը դրանցում հանդիպող լուծումների բնույթի վերաբերյալ և տալիս է նորմերի դասերի համար Ֆերմա--Տորիչելիի պրոբլեմի վիճակի ամբողջական պատկերը:
\newpage
\section*{Գրականության ցանկ}
\addcontentsline{toc}{section}{Գրականության ցանկ}

\begin{enumerate}[label={[\arabic*]}]
    \item Bajaj C. 1988, “The algebraic degree of geometric optimization problems.“, Discr. Comput. Geom., 3, 177–191.
    \item Bannikova A. G., Ilyutko D. P., Nikonov I. M. 2016, “The Length of an Extremal Network in a Normed Space: Maxwell Formula“, Journal of Mathematical Sciences, 214:5, 593–608.
    \item Boltyanski V., Martini H., Soltan V. 1999, “Geometric methods and optimization problems.“, Kluwer Acad. Publ.
    \item Brazil M., Graham R. L., Thomas D. A., Zachariasen M. 2013, “On the History of the Euclidean Steiner Tree Problem.“
    \item Cieslik D. 1988, “The Fermat-Steiner-Weber-problem in Minkowski spaces.“, Optimization 19, 485–489.
    \item Cockayne E. J., Melzak Z. A. 1969, “Euclidean constructibility in graph-minimization problems“, Math. Mag., 42, 206–208.
    \item Durier R., Michelot C. 1985, “Geometrical properties of the Fermat-Weber problem.“, Europ. J. Oper. Res. 20, 332–343.
    \item Ilyukhin D. A. 2022, “The Fermat–Torricelli problem in the case of three-point sets in normed planes“, Chebyshevskii Sbornik. Vol. 23(5): 72-86.
    \item Ilyukhin D. A. 2024, “The Fermat-Torricelli problem in normed spaces“, Facta Universitatis (Nis), Ser. Math. Inform. Vol. 39(5): 943–958.
    \item Ivanov A. O., Tuzhilin A. A. 2002, “Branching geodesics in normed spaces“, Izvestiya: Mathematics, 2002, 66:5, 905–948.
    \item Jarník V., Kössler M. 1934, “O minimálních grafech obsahujících n daných bodů“, Čas. Pěstování Mat. (Essen) Т. 63: 223—235.
    \item Kupitz Y. S., Martini H. 1997, “Geometric aspects of the generalized Fermat–Torricelli problem.“, Bolyai Society Mathematical Studies 6, 55-127.
    \item Martini H., Swanepoel K. J., Weis G. 2002, “The Fermat–Torricelli problem in normed planes and spaces.“, Journal of Optimization Theory and Applications 115, 283-314.
    \item Nguyen S. D. 2018, “Constrained Fermat-Torricelli-Weber Problem in real Hilbert Spaces“, ArXiv e-prints, arXiv:1806.04296.
    \item Torricelli E. 1919, “De maximis et minimis.“ Opere di Evangelista Torricelli, Faenza, Italy.
    \item Uteshev A. Y. 2014, “Analytical solution for the generalized Fermat–Torricelli problem“, The American Mathematical Monthly 121(4), 318–331.
    \item Zachos A. N. 2014, “An analytical solution of the weighted Fermat-Torricelli problem on the unit sphere“, ArXiv e-prints, arXiv:1408.6495.
\end{enumerate}


\end{document}